\documentclass{article}
\usepackage{graphicx} % Required for inserting images
\usepackage{James}
\usepackage{dynkin-diagrams}
\usepackage{bussproofs}
\title{Model Theory and Non-classical Logic(s)}
\author{James Hindmarch}
\setcounter{section}{-1}
\numberwithin{thm}{subsection}
\date{October 2023}
\begin{document}

\maketitle
\tableofcontents
\newpage
\newlec
\section{Preliminaries}
Notation: If $\calL$ is a first order language, and $\calM$ is an $\calL$-structure and $\calA \subseteq \calM$ is a subset, we'll write $\calL_{\calA}$ for the languagte obtained by adding a new constant synbol to the signature of $\calL$ for each $a \in \calA$. $\calM$ is then naturally also an $\calL_{\calA}$ structure. \gap 
Convention: $\emptyset$ is considered an $\calL$-structure in this course. (Not everyone adopts this convention).
\begin{displayquote}
    Itaque illos planos, quos mathematicos vocant, plane consulere non desistebam, quod quasi nullum eis esset sacrificium, et nullae preces ad aliquem spiritum ob divinationem dirigerentur. quod tamen Christiana et vera pietas consequenter repellit et damnat.
\end{displayquote}
- Saint Augustine of Hippo (English below)
\begin{displayquote}
    Those impostors then, whom they style Mathematicians, I consulted without scruple; because they seemed to use no sacrifice, nor to pray to any spirit for their divinations: which art, however, Christian and true piety consistently rejects and condemns
\end{displayquote}
\section{Model Theory}
\subsection{Substructures and diagrams}
\begin{fancydefine}{$\calL$-homomorphism}
    Let $\calM$ and $\calN$ be $\calL$-structures. A $\calL$-homomorphism is a map \[
    \eta : \calM \to \calN \]
    such that, given $\bar{a}=(a_{1},\dots,a_{n}) \in \calM^n $, for all function symbols of arity $n$, we have \[
    \eta(f^{\calM}(\bar{a})) = f^{\calN}(\eta(\bar{a})). \] And for all relation symbols $\calR$ of arity $n$ we have \[ \bar{a} \in \calR^{\calM} \Leftrightarrow \eta(\bar{a}) \in \calR^{\calN}.\]
    We will call \textit{injective $\calL$-homomorphisms} $\calL$\textbf{-embeddings.} We will call \textit{invertible $\calL$-homomorphisms} $\calL$\textbf{-isomorphisms.}
    \gap
    If $\calM \subseteq \calN$ and the inclusion map is an $\calL$-homomorphism, we say that $\calM$ is a \textbf{substructure} of $\calN$ and that $\calN$ is an \textbf{extension} of $\calM$.
\end{fancydefine}
\ex \hspace{0.1cm}
\begin{enumerate}
    \item $\calL$ the language of groups (without the inverse operation, since it isn't strictly necessary), then $(\bbN, \, +, \, 0)$ is a substructure of the integers with the same operations, but not a subgroup.
    \item If $\calM$ is an $\calL$-structure, and $X \subseteq \calM$, then $X$ is the domain of a substructure of $\calM$ if and only if it is closed under the interpretation of all the function symbols (exercise).
\end{enumerate}
The \textbf{substructure generated by $X$} (i.e. the smallest substructure of $\calM$ which includes $X$) is the smallest subset of $\calM$ which contains $X$ and is closed under the interpretation of function symbols (also called the Skolem Hull). \gap
We denote the substructure generated by $X$ as $\langle X \rangle_{\calM}$. It is easy to check that $|\langle X \rangle_{\calM}| \leq |X| + |\calL|$. (Where $\calL$ denotes the cardinality of the number of sentences in $\calL$, and not the cardinality of the signature of $\calL$.) \gap
We notice a fact:
\[ (\bbR, \, \times, \, -1) \models \neg \exists x (x^{2} = -1) \]
but
\[ (\bbC, \, \times \, -1) \notmodels \neg \exists x (x^{2} = -1). \]
However $\bbR$ is a substructure of $\bbC$ (in some unnecessary to specify language), So what sorts of sentence \textit{are} preserved under substructures? \gap
Quantifier-free sentences are preserved! "They lack context."
\begin{fancyprop}{}
Let $\phi(\bar{X})$ be a quantifier-free $\calL$-formula with $n$ variables, $\calM$ an $\calL$-structure, and $\bar{a} \in \calM^{n}$. \gap
Then for every extension $\calN$ of $\calM$, we have that $\calM \models \phi(\bar{a}) \Leftrightarrow \calN \models \phi(\bar{a})$.
\end{fancyprop}
\proof We proceed by induction on the structure of formulae. \gap
First, we show that if there is some term $t(\bar{x})$ with $k$ free variables then \[ t^{\calM}(\bar{b}) = t^{\calN}(\bar{b}) \] whenever $\bar{b} \in \calM^{k}$. This is clearly the case if $t=x_{i}$ is a variable as both structures interpret $t(\bar{b})$ as $b_{i}$. \gap
If $t$ is of the form $f(q_{1}, \dots, q_{l})$ for some function $f$ of arity $l$, and some terms $q_{i}$, then by induction we have that the $q_{i}$'s coincide for both structures, i.e. $q_{i}^{\calM}(\bar{b}) = q_{i}^{\calN}(\bar{b})$; therefore
\begin{align*}
    t^{\calM}&=f^{\calM}(q_{1}^{\calM}(\bar{b}) \dots q_{l}(\bar{b}))\\
    &= f^{\calN}(q_{1}^{\calM}(\bar{b}) \dots q_{l}^{\calM}(\bar{b}))\\
    &= f^{\calN}(q_{1}^{\calN}(\bar{b}) \dots q_{l}^{\calN}(\bar{b}))\\
    &= t^{\calN}(\bar{b})
\end{align*}
For $t_{1},\,t_{2}$ some free variables, we have
\begin{align*}
    \left(\calM \models (t_{1}(\bar{a}) = t_{2}(\bar{a}))\right) &\Leftrightarrow t_{1}^{\calM}(\bar{a})=t_{2}^{\calM}(\bar{a})\\
    &\Leftrightarrow t_{1}^{\calN}(\bar{a}) = t_{2}^{\calN}(\bar{a})\\
    &\Leftrightarrow \left(\calN \models (t_{1}(\bar{a}) = t_{2}(\bar{a}))\right).
\end{align*}
Then for other atomic formulae, we let
\[ \calR = (t_{1}, \dots, t_{n}) \]
be an $n$-arity relation symbol. Then
\[ \calM \models \calR(t_{1}(\bar{a}), \dots , t_{n}(\bar{a})) \]
if and only if
\[ t_{1}^{\calM}(\bar{a}),\dots,t_{n}^{\calM}(\bar{a}) \in \calR^{\calM}\]
which happens precisely if \[ (t_{1}^{\calM}(\bar{a}),\dots, t_{n}^{\calM}(\bar{a})) \in \calR^{n} \]
as $\calN$ extends $\calM$. But then the interpretation coincides, so
\[ (t_{1}^{\calN}(\bar{a}), \dots, t_{n}^{\calN}(\bar{a})) \in \calR^{\calN}\]
\[ \text{i.e. } \calN \models \calR(t_{1}(\bar{a}),\dots, t_{n}(\bar{a})).\]
Now we just need to clean up $\neg$ and $\wedge$: \gap
$\calM \models \neg \phi$ iff $\calM \notmodels \phi$ iff $\calN \notmodels \phi$ iff $\calN \models \neg \phi$.
\gap
and $\calM \models \phi \wedge \psi$ iff $\calM \models \phi$ and $\calM \models \psi$ iff $\calN \models \phi$ and $\calN \models \psi$ iff $\calN \models \phi \wedge \psi$. Now we can build everything quantifier free. \flushright $\Box$ \flushleft 







\newlec
\begin{fancydefine}{Elementarily equivalent/Elementary embedding}
    Structures $\calM$ and $\calN$ are \textbf{elementarily equivalent} if for every $\calL$-sentence
    \[ \calM \models \phi \Leftrightarrow \calN \models \phi. \]
    A map $f: \calM \to \calN$ is an \textbf{elementary embedding} if:
    \begin{itemize}
        \item It is injective
        \item $\forall \calL$-formulae $\phi(x_{1}, \dots, x_{n})$ and elements $m_{1}, \dots, m_{n} \in \calM$, we have that
        \[ \calM \models \phi(m_{1}, \dots, m_{n}) \Leftrightarrow \calN \models \phi(f(m_{1}), \dots, f(m_{n})).\]
    \end{itemize}
    This implies $\calM$ and $\calN$ are elementarily equivalent, and is therefore stronger.
    If $\calM$ and $\calN$ are elementarily equivalent, then we write \[ \calM \equiv \calN.\]
\end{fancydefine}
\remark If $\calM$ and $\calN$ are $\calL$-structures and $\bar{m} \in \calM$, $\bar{n} \in \calN$ are ordered tuples of the same size, then \[
(\calM, \bar{m}) \equiv (\calN, \bar{n})\]
means that these are equivalent when considered as $\calL$-structures with $k$ many constants: interpreted as $\bar{m}$ in $\calM$ and as $\bar{n}$ in $\calN$. That is to say, they are equivalent as $\calL_{\bar{c}}$-structures
\begin{fancyprop}{} If $\calM \cong \calN$, then $\calM \equiv \calN$.\end{fancyprop}
\proof Left as exercise, but proceed by induction over the structure of formulae. \gap
Recall that a theory $\calT$ is complete if $\calT \proves \varphi$ or $\calT \proves \neg \varphi$ for every sentence $\varphi$. Any two models of the same complete theory are elementarily equivalent, but the models can have different cardinalities (and therefore not be isomorphic, so the implication can only go one way).
\begin{fancydefine}{Elementary substructure} A substructure $\calM \subseteq \calN$ is an \textbf{elementary substructure} if the inclusion map \[
\iota : \calM \to \calN \]
is an elementary embedding.
\end{fancydefine}
\begin{fancydefine}{Model-complete} A theory $\calT$ is called \textbf{model-complete} if every embedding between models of $\calT$ is elementary. (For example, the theory of algebraically closed fields is model-complete).
\end{fancydefine}
\begin{fancydefine}{$\kappa$-categorical} Let $\kappa$ be an infinite cardinal. We say that a theory $\calT$ is $\kappa$-categorical if all models of $\calT$ of cardinality $\kappa$ are isomorphic.
\end{fancydefine}
    
\begin{fancyprop}{Vaught's Test}
Let $\calT$ be a consistent theory with no finite models. If $\calT$ is $\kappa$-categorical for some infinite $\kappa \geq |\calL|$ then $\calT$ is a complete theory.
\end{fancyprop}
\proof Suppose $\calT$ not complete and take $\varphi$ such that \[
\calT \nvdash \varphi \text{  and  } \calT \nvdash \neg \varphi\]
Then $\calT \cup \{\varphi\}$ and $\calT \cup \{\neg\varphi\}$ are consistent, so they have infinite models (since $\calT$ does). Then by Upward L\"owenheim-Skolem, they have models of size $\kappa$. As they are models of $\calT$ of size $\kappa$, they must be isomorphic. Howeve one model believes $\varphi$ and one believes $\neg\varphi$.\\ \flushright $\contradiction.$ $ \Box$ \flushleft
\exs \hspace{0.1cm}
\begin{enumerate}
    \item Any two countable dense linear orders without endpoints must be isomorphic to $\bbQ$, so Dense Linear Orders without endpoints are $\aleph_{0}$-categorical, hence complete.
    \item For any field, $F$, the theory of infinite vector spaces is $\kappa$-categorical for $\kappa> |F| $ (exercise). Hence this theory is complete.
\end{enumerate}

\begin{fancyprop}{Tarski-Vaught Test} (\textit{not} the same as Vaught's test)\\
Suppose $\calN$ is an $\calL$-structure and $M \subseteq \calN$, then $M$ is the domain of an elementary substructure if and only if, whenever $\bar{m} \in M$ and $\varphi(x,\bar{t})$ is a formula:
\[ \exists n \in \calN \text{  s.t.  } \varphi(n,\bar{m}) \text{ \normalfont implies that  } \exists \hat{m} \in M \text{  s.t.  } \calN \models \varphi(\hat{m}, \bar{m}).\]
\end{fancyprop}
\proof If $M$ is the domain of an elementary substructure, then this is clear, since \[\calN \models \exists x. \, \varphi(x,\bar{m}) \, \, \Rightarrow \, \, \calM \models \exists x. \, \varphi(x,\bar{m}).\] So we have that \[\calM \models \varphi(\hat{m},\bar{m})\] for some $\hat{m} \in M$. But then \[
\calN \models \varphi(\hat{m},\bar{m}).\]
Conversely, if $M \subseteq \calN$ has the stated property, consider the formulae:
\[ \varphi_{f}(x,\bar{t}) := x = f(\bar{t}) \] for each function symbol $f$ in $\calL$. So for any $\bar{m} \in M$, there is an $n \in \calN$ such that \[
\calN \models (n=f(\bar{m})),\]
so by hypothesis, there is $\hat{m} \in M$ such that $\calN \models (\hat{m} = f(\bar{m}))$. So it must be the case that $M$ is closed under the interpretations of function symbols. \gap
Interpret relation symbols via \[ \calR^{\calM} = \calR^{\calN} \cap \calM^{k},\]
where $k$ is the number of elements in the tuples of the relation. Through this interpretaion, we turn $M$ into an $\calL$-structure which is clearly a substructure of $\calN$.
That $M$ is an \textit{elementary} substructure is clear, since quantifier free sentences must already hold by proposition 1.1.3, but then we can induct over formulae that contain existence quantifies by out hypothesis. \flushright $\Box$ \flushleft
\begin{fancydefine}{Universal Formula} A \textbf{universal formula} is one of the form \[ \forall \bar{x}. \, \varphi(\bar{x},\bar{y}),\] where $\varphi$ is quantifier-free.
\end{fancydefine} 
\begin{fancydefine}{Universal Theory} A \textbf{universal theory} is one whose axioms are (or can be) \textbf{universal sentences}.
\end{fancydefine}
\begin{fancydefine}{Diagram of an $\calL$-structure}
    Let $\calN$ be a $\calL$-structure, we define the \textbf{diagram} of $\calN$ to be: \[\diag(\calN) = \{ \varphi(n_{1}, \dots, n_{k}) \, : \, \varphi \text{ is a quantifier-free $\calL_{\calN}$ formula,}\] \[\text{ and $\calN \models \varphi(n_{1},\dots,n_{K})$}\}.\]
    The \textbf{elementary diagram} of $\calN$ is:
    \[\eldiag(\calN) = \{ \varphi(n_{1}, \dots, n_{k}) \, : \, \varphi \text{ is a $\calL_{\calN}$ formula, and $\calN \models \varphi(n_{1},\dots,n_{K})$}\}.\]
\end{fancydefine}
A model of the diagram of $\calN$ is an extension of $\calN$, and a model of the elementary diagram of $\calN$ is an elementary extension of $\calN$.
\begin{fancylem}{} Let $\calT$ be a consistent theory, and let $\calT_{\forall}$ be the theory of all universal sentences which $\calT$ proves. If we take a model $\calN$ of $\calT_{\forall}$, then $\calT \cup \diag(\calN)$ is consistent.
\end{fancylem}
\proof Suppose $\calT \cup \diag(\calN)$is not consistent. We know that $\calT$ is consistent itself. Then compactness tells us that there is then a finite number of sentences in $\diag(\calN)$ which are inconsistent with $\calT$ (i.e. the ones that show up in the proof of $\bot$). \gap If $\varphi(\bar n)$ is their conjunction, then $\calT \cup \{ \varphi(\bar n) \}$ is inconsistent, that is: \[ \calT \proves \neg \phi(\bar n).\]
But $\calT$ has nothing to say about $\bar n$, since $\calT$ is an $\calL$ theory, not a $\calL_{\calN}$ theory (i.e. it does not 'recognise the existence of $\bar n$ specifically). Therefore \[ \calT \proves \forall \bar x.\, \neg \varphi (\bar x).\] This is a universal consequence of $\calT$, so it must hold in $\calN$. \endproof
\begin{fancythm}{Tarski, L\=os}
An $\calL$-theory $\calT$ has a universal axiomatization if and only if whenever $\calN$ is a substructure of $\calM$ and $\calM \models \calT$, then $\calN \models \calT$.
\end{fancythm}
\proof One direction is already done by proposition 1.1.3. So we proceed with the converse. \gap Suppose that $\calT$ is preserved under substructures. Then if $\calN \models \calT$, naturally $\calN \models \calT_{\forall}$. \gap Conversely, if $\calN \models \calT_{\forall}$, then by the previous lemma 1.1.13, \[ T \cup \diag(\calN)\] is consistent. Let $\calN^{*}$ be a model of $\calT \cup \diag(\calN).$ Then $\calN^{*}$ is an extension of $\calN$ and models $\calT$. So \[ \calN \subseteq \calN^{*} \models \calT. \]
Now, by the hypothesis on $\calT$, it is preserved under substructures, so \[ \calN \models \calT.\] Therefore $\calT$ and $\calT_{\forall}$ have the same models.
\endproof
We can obtain yet more results with the same 'method of diagrams.' \gap
\textbf{1. Finding a common elementary extension of given structures:} \gap
\begin{fancythm}{Elementary amalgamation}
    Let $\calM$ and $\calN$ be $\calL$-structures, and let $\bar m \in \calM^{k}$, $\bar n \in \calN^{k}$ such that \[ (\calM,\bar m) \equiv (\calN, \bar n).\] Then there is an elementary extension $\calK$ of $\calM$ such that there is an elementary embedding \[
        g : \calN \hookrightarrow \calK
    \]
    with $g(\bar n) = \bar m$.
\end{fancythm}
\proof By replacing $\calN$ with an isomorphic copy (if need be), we will assume $\bar m = \bar n$. \gap We'll show that \[ \calT := \eldiag(\calM) \cup \eldiag(\calN)\] is consistent (and, of course, we will use compactness). Suppose we are given a finite subset $\Phi$ of sentences in $\calT$ which includes a finite number of sentences from $\eldiag(\calN)$. Let the conjunction of these sentences be $\varphi (\bar m, \bar k)$ where $\phi(\bar x, \bar y)$ is a $\calL_{\calN}$-formula and $\bar k$ are pairwise distinct elements that occur in $\calN$ but not in $\bar m$. (That last sentence was just bookkeeping, don't think it means more than it does). \gap 
If $\Phi$ is inconsistent then \[ \eldiag(\calM) \proves \neg \varphi(\bar m, \bar k).\] Since elements in $\bar k$ are distinct and not in $\calM$, we must have: \[ \eldiag{\calM} \proves \forall \bar y. \, (\neg \varphi(\bar m, \bar k))\] So we have \[ (\calM,\bar m) \models \forall \bar y. \, \neg \varphi(\bar m, \bar y).\] But, by hypothesis,
\[ (\calM, \bar m) \equiv (\calN, \bar n).\] Therefore, it must be the case that \[ (\calN, \bar n) \models \forall \bar y \neg \varphi(\bar m, \bar y)\] But we assumed that $\varphi(\bar m, \bar k) \in \eldiag(\calN)$ $\contradiction$ By compactness, $\calT$ is consistent, then we can take $\calK$ to be the $\calL$-reduct of some model of $\calT$.
\endproof
\textbf{2. Controlling the size of a model:} \gap 
\begin{fancythm}{Stronger L\"owenheim-Skolem}
    Let $\calM$ be an infinite $\calL$-structure, and $\kappa \geq |\calL|$ be an infinite cardinal. Then \begin{itemize}
        \item[$(\downarrow)$] If $\kappa < |\calM|$, then $\calM$ admits an elementary substructure of size $\kappa$.
        \item[$(\uparrow)$] If $\kappa > |\calM|$, then $\calM$ admits an elementary extension of size $\kappa$
    \end{itemize}
\end{fancythm}
\proof \hspace{0.1cm}
 \begin{itemize}
    \item[$(\downarrow)$] We will do this later, when we have a bit more firepower.
    \item[$(\uparrow)$] Expand $\calL$ by adding a constant for everything in $\calM$, and for everything in $\kappa$. Then let $\calT$ be \[\eldiag(\calN) \cup_{i\neq j \in \kappa} \{ c_{i} \neq c_{j} \}. \]
    \item[] By compactness, $\calT$ has a model, and this model must be of size $\geq \kappa$. Then we use $(\downarrow)$ to get of size \textit{exactly} $\kappa$. (In particular, a model of a countable language has a countable elementary substructure.)  
\end{itemize}










\newlec 
If we take some theory, is there some way to \textit{make} it universally axiomatizable?
\subsection{Existentially closed structures and quantifier elimination}
\begin{fancydefine}{Skolem Function/Skolemisation}
     Say we have an $\calL$-theory $\calT$, and some formula $\varphi(\bar x, \bar y)$ with non-empty $\bar x$. A \textbf{Skolem function} for $\varphi$ is an $\calL$-term $t$, such that: \[ \calT \proves \forall \bar x \, (\exists y. \varphi(\bar x, y) \Rightarrow \phi(\bar x, t(\bar x)))\]
A \textbf{Skolemisation} of an $\calL$-theory is a language $\calL^{+} \supseteq \calL$, together with an $\calL^{+}$-theory $\calT^{+} \supseteq \calT$ such that:
\begin{enumerate}
    \item Every $\calL$-structure that models $\calT$ can be expanded to a model of $\calT^{+}$
    \item The theory $\calT^{+}$ admits Skolem functions for every $\calL^{+}$-formula $\varphi(\bar x, y )$ with $\bar x \neq \emptyset$.
\end{enumerate}
A theory $\calT$ is a \textbf{Skolem theory} if it is a skolemisation of itself.
\end{fancydefine}
\begin{fancyprop}{Elimination of quantifiers?}
    Let $\calT$ be an $\calL$-theory, and $F$ a collection of $\calL$-formulae, which includes all the atomic formulae and is closed under boolean combinations ($\vee, \neg, \wedge$). \gap Then if for every formula in $\psi(\bar x,y)$ in $F$, we have $\varphi( \bar x)$ in $F$ such that \[ \calT \proves \forall \bar x \, (\exists y \psi(\bar x, y) \Leftrightarrow \phi(\bar x)),\] then every $\calL$-formula is equivalent (according to $\calT$, we will say \textit{modulo} $\calT$) to a formula in $F$ with the same free variables 
\end{fancyprop}
\proof Example Sheet 1.
\begin{fancyprop}{}
    Let $\calT$ be a Skolem $\calL$-theory. Then:
    \begin{enumerate}
        \item Every $\calL$-formula $\varphi( \bar x)$ is equivalent to some quantifier-free $\varphi^{*}(\bar x)$ modulo $\calT$. ($\bar x \neq \emptyset$).
        \item If $\calN \models \calT$ and $X \subseteq \calN$, then either $\langle X \rangle_{\calN} = \emptyset$ or $\langle X \rangle_{\calN} \preceq \calN$.
    \end{enumerate}
\end{fancyprop}
\proof \hspace{0.1cm}
\begin{enumerate}
    \item Clearly $\varphi(\bar x, t(\bar x))$ implies $\exists y.\, \varphi(\bar x, y)$ in any model, so having Skolem functions means that \[ \calT \models \forall \bar x (\exists \varphi(\bar x, y) \Leftrightarrow \varphi(\bar x, t(\bar x))).\] So we are done by 1.2.2.
    \item If $\calM = \langle X \rangle_{\calN}$ is the Skolem hull, $\bar m \in \calM$, and $\varphi(\bar x, y)$ is such that $\calN \models \exists y. \varphi(\bar m, y)$, then there is an $\calL$-term $t$ (the associated Skolem function) such that $\calN \models \varphi(\bar m, t(\bar m))$. \gap Since $\calM$ is closed under function symbols \[ t^{\calN}(\bar m) \in \calM.\] By the Tarski-Vaught test (Prop 1.1.11), we have a witness to any existential statement, so $\calM \preceq \calN$. \endproof
\end{enumerate}
Skolem theories aren't usually found in nature (they are relatively unnatural, in fact). Fortunately, we \textit{can} construct them.
\begin{fancythm}{Skolemisation Theorem}
    Every (first-order) language $\calL$ can be expanded to some $\calL^{+} \supseteq \calL$ contianing $\Sigma$ such that:
    \begin{enumerate}
        \item $\Sigma$ is a Skolem $\calL^{+}$-theory,
        \item Any $\calL$-structure can be expanded to a $\calL^{+}$-structure that models $\Sigma$.
        \item $|\calL| = |\calL^{+}|$
    \end{enumerate}
\end{fancythm}
\idea Our language might lack Skolem functions, so we just need to add them in. We design $\calL^{+}$ to include Skolem functions for each suitable formula.
\proof If $\chi(\bar x, y)$ is an $\calL$-formula ($\bar x \neq \emptyset$), then we add a function symbol $F_{\chi}$ of arity $|\bar x|$. By doing this, we obtain a language $\calL'$ containing all these new functions. \gap Next, we define $\Sigma(\calL)$ to be the set of $\calL'$-sentences, which enforces that these function symbols are actually Skolem functions: \[ \forall \bar x \, (\exists y. (\chi(\bar x,y) \Rightarrow \chi(\bar x, F_{\chi}(\bar x)))). \tag{$*$} \]
There is a problem though, $\Sigma(\calL)$ is an $\calL'$ theory, not an $\calL$ theory. The solution is to iterate this $\omega$ times, as each stage $\Sigma_{n+1}(\calL)$ then says that you have Skolem functions for $\calL_{n}$-formulae at the previous step\footnote{We need to add Skolem functions for the new language's formulae, then add Skolem functions for those formulae, so we need a fixed point.}. \gap 
Start with $\calL_{0}= \calL$, and $\Sigma_{0} = \emptyset$. Then we define recursively: \[ \calL_{n+1} = \calL_{n}'\] and \[ \Sigma_{n+1}=\Sigma_{n} \cup \Sigma(\calL_{n}).\] Then \[ \calL^{+} = \bigcup_{n<\omega} \calL_{n}\] \[ \Sigma = \bigcup_{n< \omega} \Sigma_{n}.\] Note that every $\calL^{+}$ formula is in some $\calL_{n}$, and thus $\Sigma_{n+1} \subseteq \Sigma$ includes a sentence saying that there is a Skolem function for it. It is also clear that $|\calL^{+}| = |\calL|$. \gap We'll show the structure expansion property (2. I think) for $\calL'$ and $\Sigma(\calL)$. Then we get an interpretation for all the formulae that we add in the expansion by just 'proceeding up the chain' in exactly the same way we define $\calL^{+}$. For this we will need the axiom of choice. 
\gap First, assume the structure is non-empty, otherwise this is trivially true (since existence claims are always false). Take $\calM \neq \emptyset$ to be an $\calL-structure$. We expand this to an $\calL'$-structure: $\calM'$ in the following way: \gap Say we have $\chi( \bar x, y)$, ($\bar x \neq \emptyset$), and a tuple $\bar m \in \calM$. \begin{itemize}
    \item If there's a $b$ such that $\calM \models \chi(\bar m, b)$, then choose one such $b$, and interpret $F_{\chi}(\bar m)$ as $b$.
    \item If there is no such $b$, then just set $F_{\chi}(\bar m) = m_{1}$ (the first entry of $\bar m$, whatever, it doesn't matter).
\end{itemize}
Then by construction ($\calM' \models (*)$) where ($*$) is the skolem sentence at the start of the proof. \endproof
\begin{corr}{} Any theory $\calT$ admits a skolemisation $\calT^{+}$ in a language of the same cardinality.
\end{corr}
\proof Take $\calT^{+}$ to be $\calT \cup \Sigma$. Any model of $\calT^{+}$ is of course a model of $\Sigma$, so $\calT^{+}$ must have skolem functions. Moreover, any $\calL$-structure $\calM$ that has $\calM \models \calT$ can be expanded to one that has $\calM^{+} \models \Sigma$, in fact, so that we have $\calM \models \calT^{+}$.
\endproof 
\begin{corr}{} Let $\calM$ be an $\calL$-structure $X \subseteq \calM$ with \[|\calL| + |X| \leq \kappa \leq |\calM|.\] Then $\calM$ has an elementary substructure $\calN$ of size $\kappa$ that contains $|X|$.
\end{corr}
\proof \newlec First, we pick a subset $Y$ of $\calM$ containing $X$, and of size $\kappa$. Let $M'$ be an expansion of $\calM$ to a Skolem theory, and consider the Skolem Hull of $Y$ in $M'$, $\langle Y \rangle_{M'}$. We know that $\langle Y \rangle_{M'}$ is an elementary substructure of $M'$, since any non-empty Skolem Hull is an elementary substructure by part 2 of 1.2.3, and so if we take $\calN$ to be the $\calL$-reduct of this $\langle Y \rangle_{M'}$, then $\calN \preceq \calM$ and $X \subseteq \calN$. Now we just need to check that the sizes work out appropriately, as: \[ |\calN|\leq |Y| + |\calL^{+}| = \kappa + |\calL| = \kappa = |Y| \leq |\calN|,\] so we have $|\calN| = \kappa$. \endproof
\begin{fancydefine}{Elimination Set}
    Let $\calT$ be an $\calL$-theory. A set $F$ of $\calL$-formulae is an \textbf{elimination set} for $\calT$ if, for every $\calL$-formula $\varphi$, there is a Boolean combination $\varphi^{*}$ of formulae in $F$ such that \[ \calT \proves \varphi \Leftrightarrow \varphi^{*} \] A theory $\calT$ has quantifier-elimination if the family of all quantifier-free $\calL$-formulae forms an elimination set of $\calT$.
\end{fancydefine}
\begin{ex}
    Consider a polynomial $p(x)=x^{3}-31x^{2}+6$. To determine whether the sentence \[ \exists x.\, p(x)=0\] holds in an algebraically closed field, we can equivalently check the quantifier-free sentence: \[ (1 \neq 0) \wedge (31 \neq 0). \] since this means there must be a root. \gap Similarly, suppose I had a matrix over $\bbR$, then the sentence which says "$M$ is invertible" is equivalent to \[ \neg (\det(M)=0).\]
\end{ex}
Why care about quantifier elimination?
\begin{enumerate}
    \item Deciding whether two models of $\calT$ are elementarily equivalent is just reduced to verifying whether they satisfy the same sentences in the elimination set.
    \item In particular, completeness of $\calT$ is reducible to checking if all sentences in $F$ are deducible from $\calT$, or inconsistent with it.
    \item If the language is recursive, and the procedure for quantifier elimination is computable, then the problem of decidability for sentences in $\calT$ is reduced to the problem of decidability for sentences in $F$.
    \item The elementary embeddings $\calM \hookrightarrow \calN$ become exactly those embedding that preserve $\varphi$ and $\neg \varphi$ for $\phi \in F$. This has nice implications about our theory being model complete\footnote{which I didn't quite hear}.
    \item Figuring out which subsets are definable is much easier once you have an elimination set. Definable subsets are simply the Boolean combinations of subsets definable by formulae in $F$.
\end{enumerate}
In what follows, we use the notation $\neg F$ for the set of all the negations of things in $F$.
\begin{fancyprop}{Syntactic Quantifier Elimination}
    Suppose we have an $\calL$-theory $\calT$, and $\calF$ a family of $\calL$-formulae which includes all the atomic formulae. Suppose that for every $\calL$-formula of the form \[ \theta(\bar x) := \exists y. \, \bigwedge_{i<n} \varphi_{i}(\bar x, y) \] with $\phi_{i} \in \calF \cup \neg \calF$, that there is a Boolean combination $\theta^{*}$ of $\calF$-formulae such that: \[ \calT \proves \forall \bar x. \, (\theta(\bar x) \Leftrightarrow \theta^{*}(\bar x)\] Then $\calF$ is an elimination set for $\calT$. 
\end{fancyprop}
\proof See proposition 1.2.2.
\begin{ex}
    The theory $\calT_{\infty}$ of infinite sets in the language with empty signature (so all atomic formulae are equalities and the only terms are variables, no function or relation symbols). This has quantifier-eliimination, since by proposition 1.2.9, it is enough to eliminate the existential quantifier in formulae $\varphi(x_{0}, \dots, x_{n-1})$ of the form: \[ \exists y.\,\left(\bigwedge_{I}y=x_{i} \wedge \bigwedge_{J}(y \neq x_{j}) \wedge \bigwedge_{K^{2} (k \neq k')} x_{k} = x_{k'} \wedge \bigwedge_{L^2 (l \neq l')} x_{l} \neq x_{l'}\right). \] WLOG, we can assume that $I=\emptyset$. Then define \[ \psi(x_{0},\dots, x_{n-1}) := \left(\bigwedge_{K^{2} \setminus K \times K} x_{k} \neq x_{k'} \wedge \bigwedge_{L^{2} \setminus L \times L}\right).\] As this formula doesn't contain $y$, we have $\phi(x_{0}, \dots, x_{n-1})$ is equivalent to \[\exists y. \, \wedge_{J} y \neq x_{j} \wedge \psi(x_{0},\dots x_{n-1}) \] We then conclude that $\varphi$ and $\psi$ are equivalent mod $\calT$, since $\calT_{\infty} \proves \forall \bar x. \, (\exists y. (y \neq x_{j}))$.
\end{ex}
\hspace{0.1cm}\\
\newlec \begin{fancydefine*}{$\Rightarrow_{1}$} If $\calM$ and $\calN$ are $\calL$-structures, we write $\calM \Rightarrow_{1} \calN$ if every existential sentence modelled by $\calM$ is also modelled by $\calN$.
\end{fancydefine*}
\begin{fancythm}{Existential Amalagmation} Let $\calM$ and $\calN$ be two $\calL$-structures, $S$ a subet of $\calM$, and $f: \langle S \rangle_{\calM} \to \calN$ a homomorphism such that \[ (\calN,f(S)) \Rightarrow_{1} (\calM,S).\] Then we can find an elementary extension of $\calM$, $\calK$, and an (not necessarily elementary) embedding $\calN \to \calK$ which makes the following diagram commute \begin{center}
    \begin{tikzcd}
        & \calK & \\
    \calM \arrow[hookrightarrow]{ur}{\preceq} & &  \arrow[hookrightarrow]{ul}[swap]{g}\calN\\
     & \langle S \rangle_{M} \arrow{ur}[swap]{f} \arrow[hookrightarrow]{ul} &
    \end{tikzcd}
\end{center} 
\end{fancythm}
\proof This is similar to theorem 1.1.15, and left as an exerice.
\endproof
We say that a class $\bbK$ of $\calL$-structures has the \underline{amalgamation property} when, given $\bbK$-structures: $\calA, \calB, \calC$ and embeddings \[ \calB \xhookleftarrow{e}\calC \xhookrightarrow{f} \calA \] there is a structure $\calD$ in $\bbK$ and embeddings \[ \calB \xhookrightarrow{g} \calD \xhookleftarrow{h} \calA.\] Making the following diagram commute:
\begin{center}
    \begin{tikzcd}
        & \calD & \\
    \calB \arrow[hookrightarrow]{ur}{g} & &  \arrow[hookrightarrow]{ul}[swap]{h}\calA\\
     & \calC \arrow{ur}[swap]{f} \arrow[hookrightarrow]{ul}{e} &
    \end{tikzcd}
\end{center} 
\begin{fancydefine}{Existentially closed}
    Let $\bbK$ be a class of $\calL$-structures, and $M \in \bbK$. We say that $\calM$ is \textbf{existentially closed} in $\bbK$ if, for every existential formula $\Psi(\bar x)$ and tuple $\bar m \in \calM$, the existence of an extension $\calM \subseteq \calN$ with $\calN \models \Psi(\bar m)$ and $\calN \in \bbK$ forces $\calM \models \Psi(\bar m)$.
\end{fancydefine}
\begin{exs} \hspace{0.1cm}
    \begin{enumerate}
        \item[(a)] Every existentially closed field is algebraically closed.\gap Let $\calA$ be existentially closed, and view a non-trivial polynomial $f( \bar y)$ over it as $p(\bar a, y)$ where $p(\bar x,y)$ is a term in $\calL_{\text{ring}}$ and $\bar a \in A$ (e.g.  See $y^{2} + 2y +3$ as $p(1,2,3,y)$ where $p(x_{0},x_{1},x_{2},y)=x_{0}y^{2}+x_{1}y+(-x_{2}).$) We can replace $f$ by an irreducible factor if need be, and take the quotient $\frac{A[y]}{(f)}$, which is a field extension of $\calA$ in which $f$ has a root, i.e. \[ \frac{A[y]}{(f)} \models \exists y.\, p(\bar a, y)=0.\] As $\calA$ is existentially closed, $\calA \models \exists y.\, p(a,y) =0$, i.e. $f$ has a root in $\calA$. In fact the existentially closed fields are precisely the algebraically closed ones. If $\calA$ is algebraically closed, then you can't solve more systems of equations and inequations over $\calA$ by considering field extensions.
        \item[(b)] The existentially closed lienar orders are the dense total orders without endpoints.
        \item[(c)] The existentially closed ordered fields are the \underline{real closed fields} (i.e. the ones that are elementarily equivalent to the real numbers, or an ordered field such that all the non-negative elements are squares, and all the odd polynomials have a root).
    \end{enumerate}
\end{exs}
\begin{fancythm}{}
    Let $\bbK$ be a class of $\calL$-structures, closed under isomorphims, and such that the class of all substructures of $\bbK$ structures has the amalgamation property. Then every existential $\calL$-formula is equivalent to a quantifier-free one in all existentially closed structures in $\bbK$. \gap In particular, if $\calT$ is a theory axiomatising existentially closed structures in $\bbK$, then $\calT$ has quantifier elimination.
\end{fancythm}
\proof For $\phi(\bar x)$ an existential formula, call a pair $(\calM, \bar m)$ a \underline{witnessing pair} if $\calM$ is existentially closed in $\bbK$, and $\calM \models \phi ( \bar m )$. For each such pair, let $\theta_{(\calM,\bar m)}(\bar x)$ be the conjunction of all the literals $\psi(\bar x)$ such that $\calM \models \psi( \bar m )$, and $\chi( \bar x )$  be the disjunction of all the $\theta_{(\calM, \bar m)}(\bar x)$ where $(\calM, \bar m)$ is a witnessing pair\footnote{These statements might be infinitary but it is okay by compactness basically.}.\gap It is enough to argue that if $\calN$ is existentially closed in $\bbK$, then $\calN \models \phi(\bar n)$ iff $\calN \models \chi(\bar n)$. If $\bar n \in \calN$ is such that $\calN \models \phi(\bar n)$, then $(\calN, \bar n)$ is a witnessing pair, and thus $\calN \models \chi( \bar n)$. \gap Conversely, if $\calN \models \chi(\bar n)$, then there must be a witnessing pair $(\calM, \bar m)$ such that $\calN \models \theta_{(\calM, \bar m)}(\bar n)$, i.e. if $\psi(\bar x)$ is a literal, and $\calM \models \psi(\bar m)$ then we must have $\calN \models \psi( \bar n)$. There is thus an embedding \[ e: \langle \bar m \rangle_{\calM} \to \calN\] which maps $\bar m$ to $\bar n$. The amalgamation property then provides some substructure $\calC$ which embeds into $\calD$ where $\calD \in \bbK$ and both $\calM$, $\calN$ embed into $\calC$ (and therefore also into $\calD$). \\
\begin{center}
    \begin{tikzcd}
        & \calD &\\
        & \calC \arrow[hookrightarrow]{u} & \\
    \calM \arrow[hookrightarrow]{ur} \arrow[uur, "g",dashrightarrow,bend left = 20] & &  \arrow[hookrightarrow]{ul} \arrow[uul,"h",dashrightarrow,swap, bend right =20]\calN\\
     & \langle \bar{m} \rangle_{M} \arrow{ur}[swap]{e} \arrow[hookrightarrow]{ul} &
    \end{tikzcd}
\end{center}

We thus have embeddings $\calM \xhookrightarrow{g} \calD \xhookleftarrow{h} \calN$ with $g(\bar m) = (h \circ e)(\bar m) = h(\bar n).$ By potentially replacing $\calD$ with an isomorphic copy, we may assume that $h$ is an inclusion, so $g(\bar m) = \bar n$.\gap We know that $\calM \models \varphi(\bar m)$, but then $\calD \models \varphi(g(\bar m))$ as $\varphi$ is existential and $\calM \subseteq \calD$. Since $\calN$ is existentially closed in $\bbK$, $\calD \in \bbK$, and $\calN \subseteq \calD$, so we conclude that $\calN \models \varphi(g(\bar m))$, i.e. $\calN \models \phi(\bar n)$. If $\calT$ axiomatises existentially closed structures in $\bbK$, then the result follows from the completeness theorem and the syntactic criterion for quantifier elimination (proposition 1.2.9). \endproof
\newlec
\begin{exs}\hspace{0.1cm}
    \begin{enumerate}
        \item The theory ACF of algebraically closed fields has quantifier elimination. Indeed, recall that ACF axiomatises the existentially closed fields. So by Theorem 1.2.14, we just need to check that the class of all substructures of fields, (i.e. integral domains) has the amalgamation property. Field theory tells us that fields have the amalgamation property, so to show integral domains do, just take the fraction field and then use the amalgamation property for fields to extend it to integral domains.
        \item The theory DLO of dense linear orders without endpoints has quantifier elimination and has the amalagamation property. This can be proved syntactically, or by using the fact that any two linear orders embed into a poset, since by Zorn's lemma, a poset embeds into a linear order.
    \end{enumerate}
\end{exs}
\begin{fancydefine}{Inductive class}
    A class $\bbK$ of $\calL$-structures is \textbf{inductive} if it is closed under isomorphisms and also under unions of embedding chains. \gap We don't talk about models, because e.g. the class of groups without an element of infinite order is an inductive class, but that property is not first-order axiomatisable. So this is slightly more general.
\end{fancydefine}
\begin{fancythm}{}
    Let $\calM$ be a structure in an inductive class $\bbK$. Then $\calM$ is a substructure of $\calN$, $\calM \subseteq \calN$ for some $\calN$ existentially closed in $\bbK$.
\end{fancythm}
\idea Just like with the proof that every field has an algebraic closure, we add elements solving all the polynomials, and then keep adding elements for the new polynomials, and take the union overall and show that this must be algebraically closed. This is just an extension of this idea.
\proof We show that $\calM$ can be extended to some structure $\calM^{*}$ which is still in $\bbK$ with the following property: \gap For all $\bar m \in \calM$ and for every existential formula $\phi(\bar x)$ if $\phi(\bar m)$ holds in some extension of $\calM^{*}$ in $\bbK$, then it holds in $\calM^{*}$. \gap 
First we show that this would be enough to prove the theorem. Indeed, we can recurively define a chain of $\bbK$-structures by \begin{align*}
    \calM^{(0)}&=\calM\\
    \calM^{(j+1)}&=\calM^{(j)}
\end{align*}
and then take its union \[ \calN = \bigcup_{j < \omega} M^{(j)}\] This is in $\bbK$ (as $\bbK$ is inductive) and extends $\calM$. Moreover, $\calN$ is existentially closed in $\bbK$. Suppose that $\phi(\bar x)$ is an existential formula, $\bar n \in \calN$, and $\calD$ is a structure in $\bbK$ such that $\calD \models \phi( \bar n)$. As \[ \bar n \in \calN = \bigcup_{j<\omega}\calM^{(j)}\] and the $\calM^{(j)}$ form a chain, $\bar n \in \calM^{(k)}$ for some $k< \omega$. So $\calM^{(k+1)} \models \phi(\bar n)$ as $\calM^{(k+1)} = (\calM^{(k)})^{*}$ and $\calD$ also extends $M^{(k+1)}$. Since existential sentences are closed under extensions, $\calN \models \phi( \bar n)$. \gap 
To show that we can build such an $\calM^{*}$, we list all the pairs $(\phi_{\beta},\bar{m}_{\beta})_{\beta < \delta}$ where $\phi$ is an existential formula, and $\bar m \in \cal M$. We then construct a chain of $\bbK$-structures by transfinite induction. We build this by setting $\calM_{0}=\calM$, and set $\calM_{\beta+1}$ to be some $\bbK$-structure $\calD$ with $\calM_{\beta} \subseteq \calD \models \phi_{\beta}(\bar m_{\beta})$, if it exists, otherwise we set it to be $\calM_{\beta}$. Then $M_{\lambda} = \cup_{\beta < \lambda} M_{\beta}$, for $\lambda$ a limit ordinal. Then $\calM^{*}=\calM_{\delta}$.
 \gap If $\phi(\bar x)$ is existential, $\bar m \in \calM$, and $\calD \in \bbK$ is such that $\calM^{*} \subseteq \calD \models \phi(\bar m)$, then $(\phi, \bar m)$ is $(\phi_{\beta}, \bar m_{\beta})$ for some $\beta< \delta$. Then $\calM_{\beta} \subseteq \calM^{*} \subseteq \calD$, so the clause in the definiton of $M_{\beta +1}$ holds true, and we must have that $\calM_{\beta +1} \models \phi(\bar m)$. As $\phi$ is existential, and $\calM^{*}$ extends $\calM_{\beta +1}$, we have $\calM^{*} \models \phi(\bar m)$. \endproof
 \begin{fancythm}{}
    Let $\calT$ be an $\calL$-theory. The following are equivalent: \begin{enumerate}
        \item[(a)] $\calT$ has quantifier elimination 
        \item[(b)] If $\calA$,$\calB \models \calT$ and $\bar a \in \calA^{n}$ and $\bar b \in \calB^{n}$ are such that $(\calA, \bar a)$ and $(\calB, \bar b)$ satisfy the same quantifier free sentences $((\calA,\bar a) \equiv_{0} (\calB), \bar b)$, then $(\calA, \bar a) \Rightarrow_{1} (\calB, \bar b)$.
        \item[(c)] Whenever $\calA, \calB \models \calT$, $S$ is a subset of $\calA$ and $e: \langle S \rangle_{\calA} \hookrightarrow \calB$, then there is an elementary extension $\calD$ of $\calB$ and an embedding $f: A \hookrightarrow \calD$ extending $e$.
        \item[(d)] The theory $\calT$ is model-complete, and $\calT_{\forall}$ has the amalgamation property.
        \item[(e)] For every quantifier free formula $\phi(\bar x,y)$ the formula $\exists y. \phi(\bar x, y)$ is $\calT$-equivalent to a quantifier-free formula $\Psi(x)$.
    \end{enumerate}
 \end{fancythm}
 \proof \hspace{0.1cm} \begin{enumerate}
    \item[$(a) \Rightarrow (b)$:] This is obvious, since if $\calT$ has quantifier elimination, then every formula is equivalent to a quantifier free formula, so $\calA$ and $\calB$ satisfy the same sentences by assumption
    \item[$(b) \Rightarrow (c)$:] It is enough to show that $(\calA, \bar a) \Rightarrow_{1} (\calB, e(\bar a))$ by theorem 1.2.11. Since a sentence in $\calL_{S}$ can only mention finitely many of the new constants in $S$, it's enough to check that $(\calA, \bar a) \Rightarrow_{1} (\calB, e(\bar a))$ for all tuples $\bar a$ obtainable from $S$. Now, if $\bar a$ is such a tuple, and $e: \langle S \rangle_{\calA} \hookrightarrow \calB$ is an embedding, then $(\calA, \bar a) \equiv_{0} (\calB, e(\bar a))$\footnote{$\equiv_{0}$ means they satisfy the same quantifier-free sentences}, so we're done by (b).
    \item[$(c) \Rightarrow (d)$:] We'll use example sheet 1 question 11(c). Fix an embedding \[ h : \calM \hookrightarrow \calN\] between models of $\calT$. Need to show that there's $\calM \preceq \calD$ and an embedding $g: \calN \to \calD$ such that $(g \circ h)(m)=m$ for all $m$. Consider the instance of (c) where $S = h(\calM) \subseteq \calN$ and $e := h^{-1} : h(\calM) \xrightarrow{\simeq} \calM$. In this case, we get an elementary extension $\calD$ with $\calM \preceq \calD$ and an embedding $g: \calN \hookrightarrow \calD$ such that $g$ extends $e$. We get \begin{center}
        \begin{tikzcd}[column sep = tiny] %Column sep is small, bottom left arrow is hook arrow for embedding
            & \calD & \\
            \calN \arrow[ur, "g"] & & \calM \arrow[ul, "\preceq",swap]\\
            & \langle S \rangle_{N} = h(\calM)\arrow[ul, hook] \arrow[ur,"e"] &\\
        \end{tikzcd} 
    \end{center}
    This means that for all $m \in \calM \preceq \calD$, we have $g(h(m))=e(h(m))=m$, as needed, so $\calT$ is model-complete. \gap To see that $T_{\forall}$ has the amalgamation property, let $\calB' \leftarrow \calC \rightarrow \calA'$ be embedding between models of $\calT_{\forall}$. By lemma 1.1.13, $\calA'$ and $\calB'$ (hence $\calC$) are substructures of models $\calA$ and $\calB$ respectively. Taking the instance of (c) where $S := \calC = \langle \calC \rangle_{\calA}$ and $e$ is $\calC \hookrightarrow \calB$ provides an elementary extension $\calD$ of $\calB$ and an embedding $f: \calA \hookrightarrow \calD$ that extends $e$. As $\calD \equiv \calB \models \calT \proves \calT_{\forall}$, the theory has the amalgamation property.
    \item[$(d) \Rightarrow (e)$] Model-completeness implies that every model of $\calT$ is existentially closed in $\Mod(\calT)$\footnote{$\Mod(\calT)$ = \{\text{ Models of $\calT$} \} } (See ES1 Q11(b)). We're then done by Theorem 1.2.14 since the models of $\calT_{\forall}$ are precisely the substructures of models of $\calT$.
    \item[$(e) \Rightarrow (a)$] This follows immediately from the syntactic criterion for quantifier elimination (proposition 1.2.9).
 \end{enumerate}
"That wasn't so bad, was it?"
\endproof
\begin{corr}{}
    Let $\calA$ be a finite $\calL$-structure, then $\Th(\calA)$ has quantifier elimination if and only if every isomorphism between finitely generated substructures of $\calA$ can be extended to an automorphism of $\calA$. \gap $\Th(\calA)$ is probably the set of sentences which hold in $\calA$.
\end{corr}
\proof \hspace{0.1cm} \begin{enumerate}
    \item[$(\Rightarrow)$] Take $\calA = \calB$ in Theorem 1.2.18(c) and \[e: \langle \bar a \rangle_{A} \xrightarrow{\simeq} \langle \bar b \rangle_{\calA} \hookrightarrow \calA. \] Then consider the elementary extension $\calD$ of $\calA$ that we obtain. \gap If $|\calA|=n<\aleph_{0}$, then $\Th(\calA)$ must include a sentence which says this, and $\calD$ also has to satisfy this sentence. We thus have that $\calA$ and $\calD$ are elementarily equivalent finite structures, so the elementary emebedding $h: \calA \hookrightarrow \calD$ must be an isomorphism. So we have \begin{center}
        \begin{tikzcd}
            \calA \arrow[r,hook, "f" ] & \calD \arrow[r, "h^{-1}"] & \calA\\
            \langle \bar a \rangle_{\calA} \arrow[u,hook] \arrow[rr,"\sim"] & & \langle \bar b \rangle_{\calA} \arrow[u,hook]
        \end{tikzcd}
    \end{center}
    So we just need to show that $f$ is also surjective. But as $|\calA| = |\calD| = n$ and $f: \calA \hookrightarrow \calD$ is an embedding, it \textit{must} also be surjective by the pigeonhole principle. So \[\calA \xrightarrow[f]{\sim} \calD \xhookrightarrow{h^{-1}} \calA\] is an automorphism of $\calA$ extending $\langle \bar a \rangle_{\calA} \xrightarrow{\sim} \langle \bar b \rangle_{\calA}$.
    \item[$(\Leftarrow)$] We prove 1.2.18(b). Let $\bar b \in \calB \models \Th(\calA)$ be as in the hypothesis for (b). As models for the complete theory $\Th(\calA)$, both $\calB$ and $\calC$ are elementarily equivalent to $\calA$. Since they're finite, they must be isomorphic WLOG. Say $\calB = \calC = \calA$. Then from \[ (\calA, \bar b) \cong_{0} (\calA, \bar c),\] we get an isomorphism $\langle \bar b \rangle_{\calA} \equiv \langle \bar c \rangle_{\calA}$ mapping $\bar b$ to $\bar c$, which we can extend to an automorphism for $\calA$. If $\bar m$ is a witness to \[ (\calA, \bar b) \models \exists \bar y. \, \varphi(\bar b, \bar y).\] then $f(\bar m)$ witnesses \[ (\calA, \bar c) \models \exists y. \, \varphi(\bar c, \bar y).\] I.e. $(\calA, \bar b) \Rightarrow_{1} (\calA, \bar c)$.

\end{enumerate}
\endproof
\begin{ex}
    Let $V$ be a finite vector space (not finite-dimensional, \textit{finite}). Any isomorphism between subspaces can be extended to an automorphism of $V$ via Steinitz, so $\Th(V)$ must have quantifier elimination.
\end{ex}
\begin{corr}{}
    Let $\calT$ be an $\calL$-theory such that:
    \begin{enumerate}
        \item Whenever I have models $\calA,\calB \models \calT$ with $\calA \subseteq \calB$, and $\varphi(\bar x, y)$ is a quantifier-free formula, and $\bar a \in \calA$ is such that $\calB \models \exists y.\, \varphi(\bar a, y)$, then $\calA \models \exists y. \, \varphi(\bar a, y)$.
        \item For any $\calC \subseteq \calA \models \calT$ both modelling $\calT$, there's an initial intermediate model $\calA' \models \calT$. That is, $\calC \subseteq \calA' \subseteq \calA$, and if $\calB$ is any other model of $\calT$ with $\calC \subseteq \calB$, then there's an embedding $\calA' \hookrightarrow \calB$ that fixes $\calC$.
    \end{enumerate}
    Then $\calT$ has quantifier elimination.
\end{corr}
\proof \newlec Let $\calA, \calB \models T$ and $\bar a \in \calA, \bar b \in \calB$ such that \[ (\calA, \bar a) \equiv_{0} (\calB, \bar b). \] It is enough to show that $(\calA, \bar a) \Rightarrow_{1} (\calB, \bar b)$. Let $\varphi(x,y)$ be quantifier free and such that $\calA \models \exists \bar y. \, \varphi(\bar a, \bar y)$. Then let $\bar c = (c_{0}, \dots, c_{k})$ be a witness to this statement, so that \[ \calA \models \varphi(\bar a, \bar c).\] 
Our strategy will be to show that we can add $\bar c$ to $\calA$ and $\bar d$ to $\calB$ so that $(\calA, \bar a, \bar c) \equiv_{0} (\calB, \bar b, \bar d)$. We claim that there is some elementary extension $\calB \preceq \calB_{0}$ and $d_{0} \in \calB_{0}$ such that \[(\calA, \bar a, c_{0}) \equiv_{0} (\calB, \bar b, d_{0}).\] If we can do this, we can iterate the procedure to get $\calB \preceq \calB_{0} \preceq \calB_{1} \preceq \dots \preceq \calB_{k-1}$ and elements $d_{j} \in \calB_{j}$ such that \[ (\calA, \bar a, c_{0}, \dots c_{j-1}) \equiv_{0} (\calB_{j-1}, \bar b, d_{0}, \dots, d_{j-1})\] for each $j<k$. Since $\varphi$ is quantifier free, we in particular have $\calB_{k-1} \models \varphi(\bar b, \bar d)$. As $\calB_{k-1} \models \exists y. \, \varphi(\bar b, \bar y)$ and $B_{k-1} \equiv \calB$, we have $\calB \models \exists y. \, \varphi(\bar b, \bar y)$ and are done.
To find $\calB_{0}$ and $d_{0}$, we use the hypotheses and compactness. As $(\calA, \bar a) \equiv_{0} (\calB, \bar b)$, there's an isomorphism \[ \langle \bar a \rangle_{\calA} \xrightarrow{\sim} \langle \bar b \rangle_{\calB}.\] Take $\calC= \langle \bar a \rangle_{\calA} \subseteq \calA$. By assumption (b) in the corollary, there's an initial intermediate model $\calC \subseteq \calA' \subseteq \calA$ of $\calT$, which must admit an embedding $\calA' \hookrightarrow \calB$ fixing $\calC$.
Then WLOG, let's assume $\calC$ is included by this embedding into $\calC \subseteq \calB$. \gap Write \[ \Psi := \{ \psi(\bar x, y) \, : \, \calA \models \psi(\bar a, c_{0}), \psi \text{ is quantifier free } \}. \] As $\bar a \in \calA'$ we have that $\calA' \models \exists y. \, \psi(\bar a, y)$ for all $\psi \in \Psi$. (by the hypothesis (a), since we have $\calA \models \exists y. \, \psi(\bar a, y)$). Now $\calA' \subseteq \calB$, and existential formula are preserved under extensions, so $\calB \models \exists y. \, \psi(\bar b, y)$ for all $\psi \in \Psi$. \gap We conclude that every finite subset of $\Psi$ is satisfied by some element of $\calB$. By compactness, there is some elementary extension $\calB \preceq \calB_{0}$ and $d_{0} \in \calB_{0}$ satisfying all the $\psi(\bar b, y)$. But then \[ (\calA, \bar a, c_{0}) \equiv_{0} (\calB, \bar b, d_{0}).\]
\endproof
\begin{ex}
    The theory RCF of real closed fields\footnote{satisfies same first order sentences as the reals, e.g. hyperreals, real algebraic numbers, the real numbers themselves.} (with signature $(+,\times, 0, 1, <)$) has quantifier elimination.
    \proof We'll assume the existence of real closures, and the fact that a real closed field satisfies the IVT for polynomials. 
    \gap First we show that 1.2.21(a) holds. Suppose we are given an embedding \[ \calA \subseteq \calB \] of real closed fields, some $\bar a \in \calA$ and a quantifier-free $\varphi(\bar x, y)$ such that $\calB \models \exists y. \, \varphi(\bar a, y)$. 
    \gap By considering the disjunctive normal form, we may assume that $\varphi$ is the disjunction of some conjunctions of literals. Moreover $y \neq z$ and $y \nless z$ can be written in terms of $=$ and $<$, so we can assume $\varphi$ is of the form: \[ \left( \bigwedge_{i < n} p_{i}(y)=0 \right) \vee \left( \bigwedge_{j<s} q_{j}>0 \right). \] where the $p_{i}$ and $q_{j}$ are polynomials with coefficients in $\calA$. 
    \gap If $\varphi$ contains some non-trivial equation $p_{i}(y)=0$, then the witness of this statement in $\calB$, if it exists, must be algebraic over $\calA$. This actually implies that the witness must already be in $\calA$ (since real closed fields are not algebraically closed, but very easy to make them algebraically closed), and we're done. Therefore we can suppose that $n=0$. 
    \gap There are only finitely many points $c_{0}, \dots, c_{n-1}$ in $\calA$ where $q_{j}=0$ for one or more $q_{j}$. Since real closed fields have the intermediate value property for polynomials, the $q_{j}$ can only change signs at one of these points. 
    \gap Note that $\calA \models \left[ (x < y) \Rightarrow \exists z. \, (x<z \vee z<y) \right]$. As the $c_{i}$ are in $\calA$, we have an element of $\calA$ between any pair of distinct $c_{i}$'s. 
    \gap Suppose that we have a witness $b \in \calB$ so that $\calB \models \varphi( \bar a, b)$. By the above, we can pick some $a \in \calA$ in the same interval between $c_{i}$ that contains $b$m and this must satisfy \[ \varphi(\bar a, a). \] So this tells us 1.2.21(a), now we check 1.2.21(b). Suppose $\calC \subseteq \calA \models \text{RCF}$. Clearly, $\calC$ must be an ordered integral domain. The field of fractions $\Frac{\calC}$ is then naturally an ordered field (say that $\frac{a}{b}>0 \Leftrightarrow ab>0$). 
    \gap The embedding of $\calC$ into $\calA$ is an injective (ordered) ring homomorphism into an (ordered) field. By the universal property of $\Frac (\calC)$, there's a unique (ordered) ring homomorphism $\Frac(\calC) \to A$ extending $\calC \hookrightarrow \calA$. Let $A'$ be the real closure of $\Frac(\calC)$. Then $\calC \subseteq \Frac( \calC) \subseteq \calA' \subseteq \calA$. If $\calB \models \text{RCF}$ and $\calC \subseteq \calB$, then by the same argument, we have a unique (ordered) ring homomorphism $\Frac (\calC) \to \calB$ extending the embedding $\calC \subseteq \calB$. Thus $\calA' \subseteq \calB$ too, and this embedding fixes $\calC$.
\end{ex}
\begin{corr}{Hilbert's Nullstellensatz}
    Let $K$ be an algebraically closed field, and $I$ be a proper ideal of $K[x_{1}, \dots, x_{n}]$, then there's an $\bar a \in K^{n}$ such that \[ f(\bar a) = 0\] for all $f \in I$.
\end{corr}
\proof By Zorn's lemma, we may assume $I$ is maximal (extend to maximal ideal $I$ containing $I'$ for any proper ideal $I'$). Let \[ L = \frac{K[x_{1}, \dots, x_{n}]}{I} \] be its residue field, and let $\bar L$ be its alfebraic closure. Choose a finite set of generators $I=(f_{1}, \dots, f_{n})$ for $I$ (which can be done by Hilbert's basis theorem). Then \[ \bar L \models \exists \bar x. \, (f_{1}(\bar x)=0 \vee \dots \vee f_{n}(\bar x)=0).\] Take $\bar x = (x_{1} + I, x_{2} + I, \dots, x_{n} + I)$. \gap
Now, we have embeddings $K \subseteq L \subseteq \bar L$, both $K$ and $\bar L$ are algebraically closed fields, and ACF is model-complete (since it has quantifier elimination). So $K \subseteq \bar L$ is elementary and therefore \[ K \models \exists \bar x. \, (f_{1}(\bar x)=0 \vee \dots \vee f_{n}(\bar x)=0).\] Take $\bar a \in K^{n}$ to be a witness to that. \endproof
\begin{corr}{}
    Let $K \models \text{ACF}$. Then the image of a cosntructible set under a polynomial map is constructible.
\end{corr}
\proof The quantifier-free-definable subsets of $K^{n}$ are precisely the finite Boolean combinations of Zariski-closed subsets. But ACF has quantifier elimination, so these are exactly the definable subsets. Now if $X \subseteq K^{n}$ is constructible, and $p: K^{n} \to K^{m}$ is a polynomial map, then \[ p(X) = \{ y \in K^{m} \, : \, \exists x. \, p(x)=y \} \] is still definable, hence constructible.
\endproof
\subsection{Ultraproducts}
\begin{fancydefine}{}
    Let $\{\calM_{i} \, : \, i \in I \}$ be a set of $\calL$-structures. The \textbf{product}: \[ \Pi_{i \in I} \calM_{i} \] of this family is the $\calL$-structure whose carrier set is \[ \Pi_{I} \calM_{i} = \{ \alpha : I \to \cup \calM_{i} \, : \, \alpha(i) \in M_{i} \}. \] with the following interpretations \begin{itemize}
        \item An $n$-ary function symbol $f$ is interpreted as \[ f^{\Pi_{I} \calM_{i}}(\Pi_{I} \calM_{i})^{n} \to (\Pi_{I} \calM_{i}).\] given by \[(\alpha_{1}, \dots, \alpha_{n}) \mapsto \lambda i. \, (\alpha_{1}(i), \dots, \alpha_{n}(i)).\] This is lambda notation representing the function mapping $i$ to $f^{\calM_{i}}(\alpha_{1}(i), \dots, \alpha_{n}(i))$.
        \item An $n$-ary relation symbol $R$ is interpreted as the subset \[ R^{\Pi_{I} \calM_{i}} \subseteq (\Pi_{I} \calM_{i})^{n}\] given by \[ R^{\Pi_{I}\calM_{i}} := \{ (\alpha_{1},\dots, \alpha_{n} \in R \, : \, (\alpha_{1}(i), \dots, \alpha_{n}(i)) \in R^{\calM_{i}} \text{ for all $i$} \}. \] 
    \end{itemize}
\end{fancydefine}
The intuition is that we want a \say{Bill} $\varphi$ to pass if \[ \{ i \, : \, \calM_{i} \models \varphi \} \] is large.
\begin{fancydefine}{Lattice/Boolean Algebra}
    A \textbf{Lattice} is a set $L$ equipped with binary commutative and associative operations $\vee$ (join) and $\wedge$ (meet) which satisfy the absorption laws. \[ a \vee (a \wedge b) = a\] and \[ a \wedge (a \vee b) \] for all $a,b \in L$. \begin{itemize}
        \item A lattice is distributive if \[ a \wedge (b \vee c) = (a \wedge b) \vee (a \wedge c).\]
        \item A lattice is bounded if there are elements $\perp, \top \in L$ such that $a \vee \perp = a$, $a \wedge \top =a$.
        \item A lattice is complemented if it is bounded and for every $a \in L$, there is $a^{*} \in L$ such that $a \wedge a^{*} = \top$ and $a \vee a^{*} = \perp$.
    \end{itemize}
    A Boolean algebra is a complemented distributive lattice.
\end{fancydefine}
Every lattice has an ordering which is induced by the algebraic structure: \[ a\leq b \Leftarrow a \wedge b = a. \]
\begin{exs}\hspace{0.1cm}
    \begin{enumerate}
        \item For every set $I$, the power set $\calP(I)$ with $\wedge = \cap$ and $\vee = \cup$ is the prototypical example of a Boolean algebra.
        \item More generally the clopen\footnote{Both closed and open} subsets of a topological space form a Boolean algebra.
        \item For any $\calL$-structure $\calM$, we have the set \[ \{ \varphi(\calM) \, : \, \varphi(x) \in \calL_{\calB} \} \] of definable subsets of $\calM$ with parameters in $\calB$ is a Boolean algebra.
    \end{enumerate}
\end{exs}
\begin{fancydefine}{Filter}
    Let $X$ be a lattice. A \textbf{filter} $\calF$ on $X$ is a subset of $X$ with the following properties: \begin{enumerate}
        \item $\calF \neq \emptyset$.
        \item $\calF$ is a terminal segment of $X$, i.e. if $f \leq x$ and $f \in \calF$, then $x \in \calF$.
        \item $\calF$ is closed under finite meets. I.e. $x,y \in \calF \Rightarrow x \wedge y \in \calF$.
    \end{enumerate}
\end{fancydefine}











\newlec \begin{exs}
    \begin{enumerate}
        \item Given an element $j \in I$, the family $\calF_{j}$ of all the subsets of $I$ containing $j$ is a filter on $\calP(I)$. Such a filter is called a \textbf{principal filter}. Filters that are not principal are called \textbf{free}.
        \item The family of all the cofinite subsets of $I$ is a filter on $\calP(I)$, called the \textbf{Fr\'echet filter}
        \item The family of all subsets of $[0,1]$ with Lebesgue measure $1$ is a filter on $\calP([0,1])$.
    \end{enumerate}
\end{exs}
\begin{fancydefine*}{Proper filter/Ultrafilter}
    A filter $\calF$ is \textbf{proper} if $\calF \neq L$, where $L$ is the full lattice. A maximal proper filter is called an ultrafilter.
\end{fancydefine*}
Ultrafilters on $\calP(I)$ are the filters $\calF$ such that for all $\calU \subseteq I$, either $\calU \in \calF$ or $I \setminus \calU \in \calF$.
\begin{fancyprop}{Ultrafilter Principle}
    Given a set $I$, every proper filter can be extended to an ultrafilter on $\calP(I)$.
\end{fancyprop}
\proof Easy application of Zorn's lemma. \endproof
\textbf{Notation: } For $\bar \alpha \in \Pi_{i \in I} \calM_{i}$ and $\varphi(\bar x)$ an $\calL$-formula, we write: \[ [\varphi(\bar \alpha)] = \{i \in I \, : \, \calM_{i} \models \varphi(\bar \alpha(i)) \}\]
\begin{fancydefine}{Reduced Product/Ultraproduct}
    Let $I$ be a set, and $\calF$ a filter on $\calP(I)$. Let $\calM_{i}$ be a family of $\calL$-structures indexed by $I$. The \textbf{Reduced Product}, which we denote as the quotient of the product by $\calF$: \[ \frac{\Pi \calM_{i}}{\calF}\] is the quotient of $\Pi_{I} \calM_{i}$ by the equivalence relation \[ \alpha \sim \beta \Leftrightarrow [\alpha = \beta] \in \calF.\] We write $\langle \alpha \rangle$ for the equivalence class of $\alpha$ (since the usual notation is already taken). \gap If the filter we took at the start of this construction was an ultrafilter, then we say that $\frac{\Pi_{I}\calM_{i}}{\calF}$ is an \textbf{ultraproduct}. If we also have all the $\calM_{i}$ are equal, we call it an \textbf{ultrapower}. We interpret the function symbols in $\frac{\Pi_{I}\calM_{i}}{\calF}$ via: \[ f^{\frac{\Pi \calM_{i}}{\calF}}(\langle \alpha_{1} \rangle, \dots, \langle \alpha_{n} \rangle) := \langle \lambda i. \, f^{\calM_{i}}(\alpha_{1}(i), \dots, \alpha_{n}(i)) \rangle. \] For relation synbols, we set \[ (\langle \alpha_{1} \rangle, \dots, \langle \alpha_{n} \rangle) \in \calR^{\frac{\Pi \calM}{I}} \Leftrightarrow [\calR(\alpha_{1}, \dots, \alpha_{n})] \in \calF.\] note that if $\calF = \calF_{j}$ is principal, then \[ \frac{\Pi \calM_{i}}{\calF} = \calM_{j}. \]
\end{fancydefine}
\begin{fancythm}{\L o\'s Theorem}
    Let $\{ \calM_{i} \, : \, i \in I\}$ be a set of $\calL$-structures, and $\calU$ be an ultrafilter on $\calP(I)$. then for every \[ (\langle \alpha_{1} \rangle, \dots, \langle \alpha_{n} \rangle) \in \left( \frac{\Pi_{I} \calM_{i}}{\calU} \right)^{n}\] and $\calL$-formula $\varphi(x_{1}, \dots, x_{n})$, we have that: \[ \frac{\Pi_{I} \calM_{i}}{\calU} \models \varphi(\langle \alpha_{1} \rangle, \dots, \langle \alpha_{n} \rangle) \,\,\,\, \text{iff} \, \,\,\, [\varphi(\alpha_{1}, \dots, \alpha_{n})] \in \calU. \] 
    In particular, if $\calM_{i} \models \calT$ for all $i \in I$ then $\frac{\Pi \calM_{i}}{\calU} \models \calT$. 
\end{fancythm}
\proof By induction over the structure of formulae. This theorem clearly holds for the atomic formulae by the definition of the interpretations in $\frac{\Pi \calM_{i}}{\calU}$. \begin{enumerate}
    \item[Negation.] If the theorem holds for $\psi$ and $\varphi = \neg \psi$, then negating both sides of \[ \frac{\Pi \calM_{i}}{\calU} \models \psi \Leftrightarrow [\psi] \in \calU \] which is the inductive hypothesis gives us \[ \frac{\Pi \calM_{i}}{\calU} \models \neg \psi \Leftrightarrow [ \psi] \notin \calU. \] But $\calU$ is an ultrafilter, so $[\neg \psi] \in \calU$.
    

    \item[Conjunction.] If it holds for $\psi_{1}$ and $\psi_{2}$, then both \begin{align*}
        \frac{\Pi_{i}\calM_{i}}{\calU} &\models \psi_{1} \,\,\,\, \text{iff} \,\,\,\, [\psi_{1}] \in \calU\\
        \frac{\Pi_{i}\calM_{i}}{\calU} &\models \psi_{2} \,\,\,\, \text{iff} \,\,\,\, [\psi_{2}] \in \calU
    \end{align*}
    Now $\frac{\Pi \calM_{i}}{\calU} \models \psi_{1} \wedge \psi_{2}$ iff $[\psi_{1}],[\psi_{2}] \in \calU$. But this is equivalent to $[\psi_{1} \wedge \psi_{2}] \in \calU$, since if $[\psi_{1} \wedge \psi_{2}] \in \calU$, then both $[\psi_{1}]$ and $[\psi_{2}]$ are in $\calU$, since $[\psi_{1} \wedge \psi_{2}] \subseteq [\psi_{1}],[\psi_{2}],$ and $\calU$ is a filter. Conversely, if $[\psi_{1}],[\psi_{2}] \in \calU$, then $[\psi_{1}] \cap [\psi_{2}] \subseteq [\psi_{1} \wedge \psi_{2}]$, so $[\psi_{1} \wedge \psi_{2}] \in \calU$ (since $\calU$ is terminal).
    \item[Existential quantification.] We finish with the case $\exists x. \, \psi(x)$, where $x$ is free in $\psi$. We have \[\frac{\Pi \calM_{i}}{\calU} \models \exists x. \, \psi(x)\] when there is $\langle \alpha \rangle \in \frac{\Pi \calM_{i}}{\calU}$ with \[\frac{\Pi \calM_{i}}{\calU} \models \psi(\langle \alpha \rangle).\] By the inductive hypothesis, this means $[\psi(\alpha)] \in \calU$. \gap
    Suppose that $\frac{\Pi \calM_{i}}{\calU} \models \psi(\langle \alpha \rangle)$. Then $[\psi(\alpha)] \subseteq [\exists x. \, \psi(x)] \in \calU$ as $\calU$ is a filter.
    Conversely, suppose that $[\exists x.\, \psi(x)] \in \calU$. Using the axiom of choice, we can pick a witness $\alpha(i)$ to $\calM_{i} \models \exists x.\, \psi(x)$ for each $i \in [\exists x. \, \psi(x)]$, and for each $i \notin [\exists x.\, \psi(x)]$, we pick an arbitrary element of $\calM_{i}$. By doing so, we can form an element $\langle \alpha \rangle \in \frac{\Pi \calM_{i}}{\calU}$ that serves as a witness to $\frac{\Pi \calM_{i}}{\calU} \models \exists x. \, \psi(x)$. \gap 
    Note that since $\calU$ is an ultrafilter, the complement of $[\exists x. \, \psi(x)]$ is not in $\calU$ which means that the set of indices $i$ for which $\alpha(i)$ was picked arbitrarily is not in $\calU$, so makes no difference in the formation of $\langle \alpha \rangle$.
\end{enumerate}
\begin{ex}{}
    The class of torsion abelian groups is not first-order axiomatisable in the language of abelian groups, i.e. the language with signature $(+,0)$.\gap
    Let $\calU$ be a free ultrafilter on the natural numbers, and consider the following ultraproduct: \[ G := \frac{\Pi C_{n}}{\calU}\] the product of cyclic groups modulo $\calU$. \gap If $C_{i+1}=\langle g_{i} \rangle$, then $g := \langle \lambda i. \, g_{i} \rangle$, the equivalence class of the sequences $g_{1},g_{2},\dots$ This has finite order if and only if \[ [ng = 0] \in \calU\] for some $n>0$. But if we fix $n$, there are only finitely many indices $i$ where $ng_{i}=0$ (the factors of $n$). However, any free ultrafilter on $\omega$ has to contain the Fr\'echet filter. Therefore the class of torsion abelian groups cannot be first order definable.
\end{ex}
\begin{ex}{}
    Fix a free ultrafilter on $\omega$, and consider the ultrapower \[ \bbN^{\calU} := \frac{\Pi_{i<\omega} \bbN}{\calU}.\] Its elements are equivalence classes of sequences $(a_{n})_{n<\omega}$, $(b_{n})_{n<\omega}$, with $\langle (a_{n}) \rangle = \langle (b_{n}) \rangle$ if $\{ n: a_{n} = b_{n}\} \in \calU$. If $\bbN$ carries its usual structure for $\calL_{\text{arithmetic}}$, then $\bbN^{\calU}$ is a nonstandard model of PA by Theorem 1.3.8. There is an embedding from the natural numbers to this set, where we map $k$ to the equivalence class of $[k,k,k,k,\dots]$. However, there is also an infinitely large natural number, given by $[1,2,3,4,\dots]$ for example, since for any $n$, the set of indices where this equivalence class is greater than $n$ is cofinite.
\end{ex}
\begin{ex}{}
    Similarly, \[ \bbR^{\calU} := \frac{\Pi_{i< \omega} \bbR}{\calU}\] is an elementary extension of $\bbR$. This includes \say{large} numbers bigger than any standard real number, e.g. $\langle (1,2,3,\dots) \rangle$. It also includes \say{infinitesimal} numbers smaller than any standard real number, e.g. $\langle (\frac{1}{1},\frac{1}{2},\frac{1}{3}) \rangle$. This is not zero, but it is smaller than any standard real number. This forms the basis for Robinson's Nonstandard Analysis.
\end{ex}
\begin{corr}{Semantic proof of compactness}
    Let $\calT$ be a first-order thory such that every finite subset of $\calT$ has a model. Then $\calT$ has a model.
\end{corr}
\proof Assume that $\calT$ is infinite, otherwise it is trivial. Let $I$ be the set of all finite subtheories of $\calT$. For each axiom $\varphi$ in the theory, we define \[ X_{\varphi} := \{ \Delta \in I \, : \, \varphi \in \Delta\},\] and let \[ D = \{ Y \subseteq I \, : \, (\exists \varphi \in \calT)(X_{\varphi} \subseteq Y)\}.\] 
Then $D$ is a proper filter on $I$, thus can be extended to an ultrafilter by the ultrafilter principle. Using the axiom of choice, pick a model $\calM_{\Delta}$ for each $\Delta \in I$, which we can do by hypothesis. 
For any $\varphi \in \calT$, we have that the set \[X_{\varphi} \subseteq \{ \Delta \in I \, : \, \calM_{\Delta} \models \varphi \}.\] Thus $\{ \Delta \in I \, : \, \calM \models \varphi \} \in D \subseteq \calU$. \gap By \L o\'s' Theorem \[ \frac{\Pi_{\Delta \in I} \calM_{\Delta}}{\calU} \models \varphi.\] So $\frac{\Pi_{\Delta \in I}}{\calU} \models \calT$.
\endproof
\subsection{Types}
\begin{fancydefine}{Definability}
    Take an $\calL$-structure $\calM$ and let $X \subseteq \calM^{n}$ be a subset of it, and let $P \subseteq \calM$. We say that $X$ is \textbf{definable with parameters in }$\mathbf{P}$ if I can find some $\bar p \in P$ and an $\calL_{P}$ formula: $\varphi(\bar x, \bar y)$ such that:\[
        X = \varphi(\calM, \bar p) = \{ \bar m \in \calM^{n} \, : \, \calM \models \varphi(\bar m, \bar p)\}.\]
\end{fancydefine}
\begin{ex}{}
    Consider the natural numbers as a structure for $\calL = \langle +, \cdot, 0, 1 \rangle$. There is an $\calL$-formula $T(e,x,s)$ such that $\bbN \models T(e,x,s)$ if and only if the Turing machine of code $e$ halts on input $x$ in at most $s$ steps. This implies that the set of all halting programmes is definable by $\exists s. \, T(e,x,s)$ (although this is not computable).
\end{ex}
\begin{fancydefine}{Lindenbaum-Tarski Algebra}
    Let $\calT$ be a theory, and let $n \in \bbN$. We obtain an equivalence relation $\sim$ on the set $\calL(\bar x)$ of $\calL$-formulas with $n$ variables, $\bar x$, by setting \[ \varphi(\bar x) \sim \psi(\bar x) \, \, \text{iff}\,\, \calT \proves \forall \bar x. \, (\varphi(\bar x) \Leftrightarrow \psi( \bar x)).\]
    The quotient \[ \calB_{n}(\calT) := \calL(\bar x) / \sim \] becomes a Boolean algebra by setting $[\varphi] \bowtie [\psi] = [\varphi \bowtie \psi]$ for any connective $\bowtie$. We call this the \textbf{Lindenbaum-Tarski algebra} of $\calT$ on variables $\bar x$.
\end{fancydefine}
\newpage
\begin{fancydefine*}{n-types}
    Let $\calM$ be an $\calL$-structure, $A \subseteq \calM$ and consider the case where $\calT$ is the $\calL_{A}$-theory $\Th_{A}(\calM)$ of all the sentences with parameters in $A$ that are true in $\calM$. The proper filters on the Boolean algebra $\calB_{n}(\calT)$ are called \textbf{n-types} of $\calM$ over $A$.
\end{fancydefine*}












\begin{fancydefine}{$n$-type/complete $n$-type}\newlec
    Let $\calM$ be a $\calL$-structure, and $A \subseteq \calM$. A set $p$ of $\calL_{A}$ formulae, with $n$ free variables $\bar x$ is an \textbf{n-type} of $\calM$ over $A$ if $p \cup \Th_{A}(\calM)$ is satisfiable. \gap
    More generally, if $\calT$ is a theory, we say that a set of $\calL$-formulae $p$ on $n$ variables $\bar x$ is an \textbf{n-type} of $\calT$ if \[\calT \cup \{\exists \bar x. \, \bigwedge \Psi \}\] is consistent for all finite subsets $\Psi \subseteq p$. \gap A \textbf{complete n-type} is a maximal $n$-type $p$, in the sense that for all $\calL$-formulae $\varphi(\bar x)$ on variables $\bar x$, we have $\varphi(\bar x) \in p$ or $\neg \varphi(\bar x) \in p$. \gap We denote the set of complete n-types of $\calT$ by $S_{n}(\calT)$ (or $S_{n}^{\calM}(A)$ if $T = \Th_{A}(\calM)$). \gap Finally, we say that $\bar m \in \calM^{n}$ \textbf{realises} an $n$-type $p$ in $\calM$ if \[ \calM \models \varphi(\bar m)\] for all $\varphi \in p$. If no $\bar m$ realises $p$, we say $p$ is \textbf{omitted} in $\calM$.
\end{fancydefine}
\begin{ex}\hspace{0.1cm}
    \begin{enumerate}
        \item Let $\calM = (\bbQ, <)$. The formulae \[ \{n < x \, : \, n \in \bbN\}\] describe a $1$-type. Any finite subset of $p \cup \Th_{\bbN}(\bbQ)$ is satisfiable in $\bbQ$. So $p \cup \Th_{\bbN}(\bbQ)$ is satisfiable by compactness. \gap 
        Obviously in $\bbQ$ itself, this type is omitted, as no rational number satsifies all of the formulae. However, this type \textit{is} realised in some elementary extension of $(\bbQ, <)$. The realisers can be thought of as imaginary, infinitely large rational numbers.
        \item Considering $\bbR$ as a structure for the theory of ordered fields, we have a set of formulae \[ \{0<x<\frac{1}{n} \, : \, n \neq 0, n \in \bbN\}\] form a $1$-type of \say{infinitesimal real numbers.}
        \item For any $\calL$-structure $\calM$, subset $A \subseteq \calM$, and element $\bar m \in \calM$, we can form the $n$-type of all, $\calL_{A}$-formulae for which $\varphi(\bar m)$ is true in $\calM$. \[ \varphi^{\calM}(\bar m / A) := \{ \varphi( \bar x) \in \calL_{A} \, : \, \calM \models \varphi( \bar m ) \}. \]
        This is a complete $n$-type, the \textbf{type of }$\mathbf{\bar m}$ over $A$.
    \end{enumerate}
\end{ex}
\begin{fancyprop}{}
    Let $\calM$ be an $\calL$-structure, $A \subseteq \calM$, and $p$ be an $n$-type of $\calM$ over $A$. Then there is an elementary extension $\calN$ of $\calM$ such that $p$ is realised in $\calN$.
\end{fancyprop}
\proof We use the method of diagrams, and show that \[ \Gamma = p \cup \eldiag(\calM)\] is satisfiable. \gap
Let $\Delta$ be a finite subset of $\Gamma$, and let \[ \varphi := \bigwedge_{\varphi' \in \Delta \cap p} \varphi'\] \[ \psi = \bigwedge_{\psi' \in \eldiag(\calM) \cap \Delta } \psi'. \]
Satisfaction of $\Delta$ is equivalent to that of \[ \phi(\bar x, \bar a) \wedge \psi(\bar a, \bar b),\]
where $\bar a \in A, \bar b \in \calM \setminus A, \varphi \in p$, and $\calM \models \varphi(\bar a_{0}, \bar b)$.\gap As $p$ is an $n$-type, there is an $\calL_{A}$ structure, $\calN_{0}$ that satisfies $p \cup \Th_{A}(\calM)$. As $\calM \models \psi(\bar a_{0}, \bar b)$, we have that \[ \exists \bar y. \, \psi(\bar a_{0}, \bar y) \in \Th_{A}(\calM).\] So $\calN_{0} \models \varphi(\bar c, \bar a) \wedge \exists \bar y. \, \psi( \bar a_{0}, \bar y)$ for some $\bar c \in \calN^{n}$. Note that $\calN_{0}$ is a $\calL_{A}$-structure, not an $\calL_{\calM}$-structure. However, by interpreting $\bar b \in \calN_{0}$ as the witness to $\exists \bar y. \, \psi(\bar a_{0}, \bar y)$, we make $\calN_{0}$ into an $\calL_{\calM}$-structure in which $\Delta$ is satisfiable. So $\Gamma$ is satisfiable by compactness. \gap Let $\calN$ satsify $\Gamma$, so that $\calM \preceq \calN$. As $\calN$ satisfies $p$, there must be a tuple $\bar n \in \calN$ with $\calN \models \varphi(\bar n)$ for each $\varphi \in p$, i.e. $\bar n$ realises $p \in \calN$. \endproof
\begin{corr}{}
    An $n$-type $p$ of $\calM$ over $A \subseteq \calM$ is complete if and only if there is an elementary extension of $\calM$: $\calM \preceq \calN$, and $\bar a \in \calN^{n}$ such that $p = $ type of $\bar a$ over $A$.
\end{corr}
\proof It is clear that if $\calM \preceq \calN$, and $\bar a \in \calN^{n}$ then the type of $a$ over $A$ in $\calN$ $\in S_{n}^{\calN}(A)=S_{n}^{\calM}(A)$. \gap Conversely, if $p$ is a complete type, then by proposition 1.4.6, there is $\calM \preceq \calN$ and $\bar a \in \calN^{n}$ that realises $p$. As $p$ is complete, if $\varphi(\bar x)$ is an $\calL_{A}$-formula, then either $\varphi \in p$ or $\neg \varphi \in p$, but not both.\gap If $\varphi \in \tp^{\calN}(\bar a /A)$, then $\calN \models \varphi(\bar a)$, so we can't have $\neg \varphi \in p$. Thus $\varphi \in p$. And if $\varphi \in p$, then $\calN \models \varphi (\bar a)$, so $\varphi \in \tp^{\calN}(\bar a /A).$ So $p= \tp^{\calN}(\bar a /A)$.
\endproof
Let $\calM$ be an $\calL$-structure, and $A \subseteq \calM$. For each formula $\varphi(x_{1}, \dots, x_{n})$, we consider the set: \[ [[\varphi]]=\{ p \in S_{n}^{\calM}(A) : \varphi \in p\} \] of complete types which include $\varphi$. Note that \[ [[\varphi \vee \psi]] = [[\varphi]] \cup [[\psi]]\] and \[ [[\varphi \wedge \psi]] = [[\varphi]] \cap [[\psi]].\] These serve as the basic open sets for a topology on $S_{n}^{\calM}(A)$. Moreover, each $[[\varphi]]$ is the complement of another open set \[[[\varphi]] = S_{n}^{\calM} \setminus [[ \neg \varphi]].\] So the topology is generated by clopen subsets. Each $S_{n}^{\calM}(A)$ is called a \textbf{Stone space}. These are compact, totally disconnected topological space.
\begin{ex}
    Let $F \models \text{ACF}$ and let $k$ be a subfield of $F$. The complete $n$-types in $S_{n}^{F}(k)$ are determined by the prime ideals of $k[X_{1},\dots, X_{n}]$. For such a type $p$, define \[ I_{p} = \{ f \in k[X_{1},\dots X_{n}] \, : \, (f(\bar x)=0) \in p\}.\]
    Each $I_{p}$ is a prime ideal, and in fact, all the prime ideals of this ring, $k[X_{1}, \dots, X_{n}]$, arise in this way. \gap The map $p \mapsto I_{p}$ is a continuous bijection $S_{n}^{F}(k) \to \Spec(k[X_{1}, \dots, X_{n}])$ where the latter is the set of prime ideals with the Zariski topology. Also note that $|S_{n}^{F}(k)| \leq |k| + \aleph_{0}$.
\end{ex}
If $p$ is isolated in $S_{n}^{\calM}(A)$, then $\{p\} = \cup_{I}[[\varphi_i]]$, so there must be a single formula $\varphi$ with $\{p\} = [[\varphi]]$. We say that $\varphi$ \textbf{isolates} $p$.
\begin{fancydefine}{Isolate}
    Let $\calT$ be an $\calL$-theory. A formula $\varphi(x_{1}, \dots, x_{n})$ \textbf{isolates} an $n$-type $p$ if:
    \begin{enumerate}
        \item $\calT \cup \varphi$ is satisfiable, and
        \item $\calT \models \forall \bar x. \, (\varphi( \bar x) \Rightarrow \psi(\bar x))$, for all $\psi \in p$. (I.e. this is true in every model of $\calT$).
    \end{enumerate}
\end{fancydefine}
\begin{fancyprop}{}
    If $\varphi$ isolates $p$, then $p$ is realised in any model of $\calT \cup \{ \exists \bar x.\, \varphi( \bar x)\}$. In particular, if $\calT$ is a complete theory, then every isolated type is realised.
\end{fancyprop}
\proof If $\calM \models \calT$, and $\calM \models \varphi(\bar a)$, then clearly $\bar a$ realises $p$ in $\calM$. If $\calT$ is complete, then either $\calT$ believes that there is an $\bar x$ satisfying $\varphi(\bar x)$, or $\calT$ believes that $\forall \bar x. \, \neg \varphi(\bar x)$. So if $\varphi$ isolates $p$, then $T \cup \{ \varphi\}$ is satsifiable by definition, and the latter possibility cannot happen.
\begin{fancythm}{Shelah's Omitting types theorem}
    Let $\calL$ be a countable language, and $\calT$ a theory in that language (not necessarily complete), and $p$ a non-isolated $n$-type of $\calT$. Then there is a countable model $\calM \models \calT$ that omits $p$.
\end{fancythm}
\proof Let $C = \{ c_0, c_1, \dots \}$ be a countable set of new constants. We expand $\calT$ to a consistent $\calL_C$ theory $T^*$ by adding recursively defined sentences $\theta_0, \theta_1, \dots$. \gap This will be done in such a way that $\theta_t \Rightarrow \theta_s$ for $s<t$. To build the $\theta_i$, we enumerate the $n$-tuples $C^n = \{ \bar d_0, \bar d_1, \dots \}$, as well as all the $\calL_C$ sentences $\varphi_0, \varphi_1, \dots$. Start with $\theta_0 = \forall x. \, x=x$ (something trivially true), and suppose we already constructed $\theta_s$ such that $\calT \cup \{\theta_s\}$ is consistent. We set $\theta_{s+1}$ as follows: \
\begin{enumerate}
    \item[$s=2i$.] These sentences are designed to turn $C$ into (the domain of) an elementary substructure of a model of $\calT$.\gap Say $\varphi_i = \exists x. \, \psi(x)$ (i.e. $\varphi_i$ is existential), and that $\calT \models (\theta_{2i} \Rightarrow \varphi_i)$. As only finitely many constants from $C$ have been used thus far, we can find some unused $c \in C$. Let $\theta_{2i+1} := \theta_{2i} \wedge \psi(c)$. \gap
    If $\calN \models \calT \cup \{ \theta_{2i} \}$, then there is a witness $a$ to $\psi$ in $\calN$. By interpreting $c$ in $\calN$ as $a$, we then have $\calN \models \theta_{2i+1}$, so $\calT \cup \{ \theta_{2i+1}\}$ is satisfiable. \gap If, on the other hand $\varphi_i$ is not existential, or $\calT \notmodels (\theta_{2i} \Rightarrow \varphi_i)$, then $\theta_{2i+1} = \theta_{2i}$.
    \item[$s=2i+1$.] We add sentences that guarantee that $C$ omits $p$. \gap Write $\bar d_{i} = (e_1, \dots, e_n)$, remove every occurence of $e_j$ from $\theta_{2i+1}$ by replacing it with the variable $x_j$, and every occurence of the other constants in $C$ by replacing them with a fresh variable $x_c$, together with a quantifier $\exists x_c$ in front of the formula. In this way we reduce $\theta_{2i+1}$ to an $\calL$ sentence instead of an $\calL_C$ sentence.\gap We obtain, therefore, an $\calL$-formula $\psi(x_{1},\dots, x_{n})$. \gap For example, if \[\theta_{2i+1} = \forall x \exists y. \, (rx +e_{1}e_{2} = y^{2}+te^{2}),\] with $r \neq t \in C$. Then \[\psi(x_{1},x_{2}) = \exists x_{r} \exists x_{c} \forall x, \exists y (x_{r}x + x_{1}x_{2} = y^{2} + x_{t}x_{2}).\] As $p$ is not isolated, there must be some $\varphi(\bar x) \in p$ that is not implied by $\psi(\bar x)$ (else $\psi$ isolates $p$). \gap 
    Define $\theta_{2i+2} := \theta_{2i+1} \wedge \neg \phi(\bar d_{i})$. Then $\calT \cup \{ \theta_{2i+2} \}$ is consistent, since there must be some $\bar n \in \calN \models \calT$ such that $\calN \models \psi(\bar n) \wedge \neg \varphi( \bar n)$, and we can turn $\calN$ into an $\calL_{C}$-structure that models $\theta_{2i+2}$ by interpreting $\bar d_{i}$ as $\bar n$ and the constants $c \in C \setminus \{ e_1, \dots, e_n\}$ as the respective witnesses to the existential statements $\exists x_i$ within $\psi$.
\end{enumerate}
Let $\calT^{*} = \calT \cup \{\theta_{0},\theta_{1}, \dots \}$. By construction $T \cup \{\theta_{s}\}$ is consistent for all $s$, and each $\theta_{s+1}$ implies $\theta_{s}$, so $\calT^{*}$ is consistent. Moreover: \begin{itemize}
    \item If $\calM$ is a model of $\calT^{*}$, the construction of $\theta_{2i+1}$ ensures that $c$ has a witness to $\varphi_{i}$ if it is an existential statement that is true in $\calM$. It is thus an elementary substructure of $\calM$ by the Tarski-Vaught test. 
    \item If $\bar c$ is some tuple in $C \models \calT^{*}$, then $\bar c = \bar d_{i}$ for some $i$. As $C \models \theta_{2i+2}$, we have $\neg \varphi(\bar c)$ for some $\varphi \in p$. So $\bar c$ can't realise $p$ in $C$.
\end{itemize}
\endproof
\subsection{Indiscernibles}
Given a linear order $\eta$, we write \[ [\eta]^{k} = \{ \bar a \in \eta^{k} \, : \, a_{0} <^{\eta} a_{1} <^{\eta} \dots <^{\eta} a_{k-1} \}. \]
\begin{fancydefine}{$\Phi$-Indiscernible}
    Let $\calM$ be an $\calL$-structure, $\eta$ be a chain (linearly-ordered but not necessarily ordered according to the structure) of elements of $\calM$, and $\Phi$ be a set of $\calL$-formulae. We say that $\eta$ is $\mathbf{\Phi}$\textbf{-indiscernible} in the structure $\calM$ if:\[
        \calM \models \varphi(\bar a) \Leftrightarrow \varphi(\bar b), \]
    for all $\bar a, \, \bar b \in \eta $ of the right size, and $\varphi \in \Phi$. \gap 
    We simply say that $\eta$ is a sequence of \textbf{indiscernibles} if the above holds for $\Phi = \{ \, \text{all $\calL$-formulae} \, \}.$
\end{fancydefine}
\begin{ex}\hspace{0.1cm}
    \begin{enumerate}
    \item Any (linearly-ordered) set of basis elements for a vector space forms a sequence of indiscernibles. Given $\bar a, \bar b \in [ \calB ]^{k}$, there is an automorphism that maps $\bar a$ to $\bar b$, so first-order formulae cannot distinguish between these elements. 
    \item Similarly, any chain of algebraically independent elements of a field $F \models ACF$ is a sequence of indiscernibles.
    \item If $R$ is a ring, then the variables $X_{1},X_{2}, \dots X_{n}$ form a set of indiscernibles of $R[X_{1}, \dots, X_{n}]$.
    \end{enumerate}
\end{ex}
\begin{fancydefine}{Ehrenfeucht-Mostowski Functor}
    An \textbf{Ehrenfeucht-Mostowski Functor} (EM Functor) is a mapping $F$ that takes each linear order $\eta$ to an $\calL$-structure $F(\eta)$, and each order-embedding to an embedding \[ \eta \xrightarrow{g} \varepsilon \] to an embedding between the corresponding structures \[ F(\eta) \xrightarrow{F(g)} F(\varepsilon), \]
    such that: \begin{itemize}
        \item Each $\eta$ generates $F(\eta)$, i.e. $\eta \subseteq F(\eta)$ as sets, and every element of the structure $F(\eta)$ is of the form $t^{F(\eta)}(\bar a)$ where $t(x_{1}, \dots, x_{n})$ is an $\calL$-term and $\bar a \in [\eta]^{k}$.
        \item For each order embedding $\eta \xrightarrow{g} \varepsilon$, the embedding of $\calL$-structures $F(\eta) F(g) F(\varepsilon)$ extends $g$.
        \item For each pair of composable embeddings $f$ and $g$, we have \[ F(g \circ f) = F(g) \circ F(f).\] And for every linear order $\eta$, we have that $F(1_{\eta}) = 1_{F(\eta)}$.
    \end{itemize}
\end{fancydefine}
In particular, every automorphism of $\eta$ induces an automorphism of $F(\eta)$.
\begin{fancyprop}{Sliding property}
    Let $F$ be an $EM$ functor, $\eta$ and $\varepsilon$ be lienar orders, and $\bar a \in [\eta]^{k}, \bar b \in [\varepsilon]^{k}$. For every quantifier-free formula on $k$ variables, $\varphi(x_{1}, \dots, x_{k})$, we have: \[ \left[ F(\eta) \models \varphi (\bar a)\right] \Leftrightarrow \left[ F(\varepsilon) \models \varphi(\bar b) \right]. \]
\end{fancyprop}
\proof Embed both $\eta$ and $\varepsilon$ into some linear order $\rho$ in which $\bar a$ and $\bar b$ get identified. Do this through maps $f$ and $g$ respectively. Suppose that $F(\eta) \models \varphi(\bar a)$. Then as embeddings preserve quantifier-free formulae, and \[ F(f): F(\eta) \hookrightarrow F(\rho) \] extends $f$, we have that $F(\rho) \models \varphi(f(\bar a))$. \gap But then $F(\rho) \models \varphi(g(\bar b))$ and so $F(\varepsilon) \models \varphi(\bar b)$ for similar reasons. \endproof So we see that the chain $\eta \subseteq F(\eta)$ is indiscernible by quantifier-free formulae. \gap Now, suppose I am given an $\calL$-structure $\calM$ which has a linear order $\eta \subseteq \calM$ inside (as a set, not as a substructure). The theory of $\eta$ in $\calM$, which we'll write $\Th(M,\eta)$, is the set of all $\calL$-formulae $\varphi(\bar x)$ which are satisfiable in $\calM$ by every ordered tuple $\bar a := a_{0} < a_{1} < \dots < a_{k-1}$ in $\eta$. The theory $\Th(F)$ of an EM-functor $F$ is the set of all $\calL$-formulae $\varphi(\bar x)$ such that $F(\eta) \models \varphi(\bar a)$ for every linear order $\eta$ and ordered tuple $\bar a$ in $\eta$.
\begin{fancylem}{}
    Let $\eta$ be an infinite linear order, $F$ be an EM functor, and $\varphi$ be a universal sentence which is true in $F(\eta)$. Then $\varphi \in \Th(F)$.
\end{fancylem}
\proof Since $\varphi$ is universal, we can write it $\varphi = \forall x. \, \psi(\bar x)$ where $\psi$ is quantifier-free. Let $\varepsilon$ be a linear order, and $\bar a$ be a tuple in $F(\varepsilon)$. We have to show that \[ F(\varepsilon) \models \psi(\bar a).\]
We know that $\varepsilon$ generates $F(\varepsilon)$, there is a finite suborder $\varepsilon_{0}$ such that $\bar a \in F(\varepsilon_{0})$. But $\eta$ is infinite, so there is an embedding $f: \varepsilon_{0} \hookrightarrow \eta$.\gap By the assumption, $F(f)(\bar a)$ satisfies $\psi$ in $F(\eta)$, so $F(\varepsilon_{0}) \models \psi(\bar a)$, since $\psi$ is quantifier-free. Similarly, $F(\varepsilon) \models \psi(\bar a)$. \endproof
\begin{fancylem}{Stretching}
    Let $M$ be an $\calL$-structure that contains the linear order $\omega$ as a generating set, and suppose that $\omega$ is an indiscernible in $\calM$ by quantifier-free formulae. Then there is an E-M Functor $F$ such that $\calM = F(\omega)$. This functor is unique up to isomorphism. That is, if $G$ is an EM-Functor with this property, then there is an isomorphism \[ F(\eta) \xrightarrow{\sim} G(\eta)\] for each linear order $\eta$, with $\alpha|_{\eta}=1_{\eta}$.
\end{fancylem}
\proof Guided proof in Example Sheet 3.\endproof 
If $F$ is an E-M functor, and $\calT$ is a theory, then the models of $\calT$ of the form $F(\eta)$ are called the \textbf{E-M models} of $\calT$.
\begin{fancythm}{Ramsey}
    Let $X$ be a countable linear order, and $k$ and $n$ be positive integers. For every \[ f: [X]^{k} \to n\] there is an infinite subset $Y \subseteq X$ such that $f$ is constant on $Y$. \gap 
    \textbf{Exercise:} Get Ya\"el to explain why this is equivalent to the colouring statement of Ramsey.
\end{fancythm}
We'll use Ramsey's theorem to show that E-M models for Skolem theories with infinite models always exist.
\begin{fancylem}{}
    Let $F$ be an E-M functor such that $\Th(F(\omega))$ is Skolem. Then $\Th(F)$ includes one of $\varphi(\bar x)$ or $\neg \varphi(\bar x)$ for every $\calL$-formula $\varphi(\bar x).$ In particular, all the structure $F(\eta)$ are elementarily equivalent and each linear order $\eta$ is indiscernible in $F(\eta)$.
\end{fancylem}
\proof Since $\Th(F(\omega))$ is Skolem, it admits a universal axiomatisation. Moreover, every formula is equivalent to a quantifier-free formula modulo $\Th(F(\omega))$. The indiscernible part of the lemma then follows from the sliding property (1.5.4) and the other part from lemma 1.5.5.
\endproof 
\begin{fancythm}{Ehrenfeucht-Mostowski}
    Let $\calM$ be an $\calL$-structure whose theory is Skolem. If $\eta$ is an infinite linear order contained in $\calM$ (as a set), then there is an E-M functor $F$ in $\calL$ whose theory expands $\Th(\calM, \eta)$.
\end{fancythm}
\proof We want to build a theory extending $\Th(\calM, \eta)$ and whose models include an indiscernible copy of $\omega$.\gap 
Expand $\calL$ by adding $\omega$-many constants \[ C = \{ c_{i} \, : \, i< \omega\}\] and build an $\calL_{C}$-theory $\calT$ whose axioms are: \[ \varphi(\bar a) \Leftrightarrow \varphi(\bar b)\] for each $\calL$-formula $\varphi(\bar x)$ and $\bar a, \bar b \in [C]^{|\bar x |}$. \gap We also add axioms \[ \varphi(c_{0},\dots, c_{k-1})\] for each formula $\varphi(x_{0},\dots, x_{k-1})$ in $\Th(\calM, \eta)$. We'll show that $\calT$ has a model by compactness. \gap Let $\calU \subseteq \calT$ be finite and list the formulae in $\calU$ as $\varphi_{0}, \dots, \varphi_{m-1}$. Note that there is some finite $k$ such that the new constants that show up in the formulae in $\calU$ are among $c_{0}, \dots, c_{k}$. Let's say that the $\varphi_{i}$ have free variables $x_{0}, \dots, x_{k-1}$ for simplicity (add some redundant variables that don't fundamentally change the formula).\gap Now we define an equivalence relation on $[\eta]^{k}$ by deciding that \[\bar a \sim \bar b\] if \[\left(\calM \models \varphi_{j}(\bar a) \right)\Leftrightarrow \left( \calM \models \varphi_{j}(\bar b)\right)\] for all $j<m$. \gap This equivalence relation partitions $[\eta]^{k}$ into finitely many equivalence classes. By Ramsey's theorem (1.5.7), thiere must be some $\bar e := e_{0}<e_{1}<\dots<e_{2k-1}$ in $\eta$ such that any two increasing $k$-tuples extracted from it land in the same equivalence class. \gap Interpreting each $c_{j}$ in $\calM$ as $e_{j}$ (for each $j<k$), we can turn $\calM$ into an $\calL_{C}$-structure that models $\calU$. So $\calT$ has a model. \gap Say $\calN \models \calT$. We know that the new constants must be interpreted as different elements of $\calN$. Since $x_{0} \neq x_{1} \in \Th(\calM, \eta)$, $\calN$ contains a copy of $\omega$ (see $c_{i}^{\calN}$ as $i$). Then consider $\calN^{*} = \calN|_{\calL}$ and $S:= \langle \omega \rangle_{\calN^{*}}$.\gap Note that $\Th(M, \eta) \subseteq \Th(\calN^{*}, \omega)$. This implies that $\Th(\calN^{*})$ is skolem, as $\Th(\calM)$ is Skolem, and $\Th(\calM) \subseteq \Th(\calM, \eta)$. It follows that $S \preceq \calN^{*}$ (see prop 1.2.3). \gap 
From this, we have that $\Th(\calM, \eta) \subseteq \Th(S,\omega)$. Finally, these sentences make $\omega$ indiscernible in $S$ by construction, so the stretching lemma (1.5.6) provides us a unique E-M functor $F$ such that $S=F(\omega)$. We're then done by lemma 1.5.8. \endproof
\section{Non-Classical Logic}
\begin{displayquote}
    There are no non-experienced truths.
\end{displayquote}
-L.E.J. Brouwer
\subsection{Intuitionistic Logic}
$\perp$ has no proof. \gap 
To prove $\varphi \wedge \psi$ is to give a proof of $\varphi$, together with a proof of $\psi$. \gap 
To prove $\varphi \rightarrow \psi$, we need a function that converts a proof of $\varphi$ into a proof of $\psi$. (In particular, to prove $\neg \varphi$, we need to convert a proof of $\varphi$ into a contradiction).\gap 
To prove $\varphi \vee \psi$ is to prove $\varphi$ or $\psi$, which lets us know which is true. \gap 
\textbf{Fact:} The law of the excluded middle is \textit{not} valid in intuitionistic logic. \gap To prove $\exists x. \, \varphi(x)$ is to give $t$ together with a proof of $\varphi(t)$.\gap 
To prove $\forall x.\, \varphi(x)$ is to give a procedure which takes in any $t$ and returns a prove of $\varphi(t)$.
\begin{fancythm}{Diaconescu}
    The Law of the Excluded Middle can be intuitionistically deduced from the Axiom of Choice.
\end{fancythm}
\proof Let $\varphi$ be a proposition. By the axiom of separation, the following are sets (within intuitionistic set theory, we mean we have a \textit{proof} that they are sets):
\begin{align*}
    A = \{ x \in \{0,1\} \, : \, \varphi \vee (x=0) \} & & B = \{ x \in \{0,1\} \, : \, \varphi \vee (x=1) \}.
\end{align*}
As $0 \in A$ and $1 \in B$, we have that $\{A,B\}$ isa family of inhabited sets. Thus, by the axiom of choice, this admits admits a choice function: \[ f: \{A,B\} \to A \cup B.\] So $f(A) \in A$ and $f(B) \in B$ by definition, hence we hae a proof of $((f(A)) = 0) \vee \varphi) \wedge ((f(B) = 1) \vee \varphi)$, and a proof that $f(A),f(B) \in \{0,1\}$. Now $f(A) \in \{0,1\}$ means we have a proof of $(f(A)=0) \vee (f(A)=1)$, and similarly for $B$. We can thus have the following situations:
\begin{enumerate}
    \item We have a proof of $f(A)=1$, so $\varphi \vee (1=0)$ is provable, so we must have a proof of $\varphi$.
    \item We have a proof of $f(B)=0$, which similarly leads to a proof of $\varphi$.
    \item Or, we could have a proof that $f(A)=0$ and $f(B)=1$, in which case we can prove $\neg \varphi$ as follows: \gap We need to turn a proof of $\varphi$ into a contradiction. \gap 
    First, take a proof of $\varphi$. Then we can prove that $A=B$ by the axiom of extensionality (i.e. they are both $\{0,1\}$). But then, $0=f(A)=f(B)=1$ as $f$ is a function, so $f$ must map the same set to the same places.
\end{enumerate}
\endproof

\begin{itemize}
    \item We assume less than classical maths, so intuitionistic logic is more general.
    \item Notions that are classically conflated, may be different in intuitionistic logic (e.g. many ways to characterise finiteness in classical logic, and they are all equivalent, but not necessarily equivalent in intuitionistic logic). More diversity of concepts.
    \item Constructive proof s have a computable content that is often missing from classical proofs.
    \item \large \textbf{Intuitionistic logic is the internal logic of an elementary topos}. \normalsize \gap This is by far the MOST important reason to study intuitionistic logic, and if you don't understand what this means, you may as well stop reading now.
\end{itemize}
What we are studying is the BHK interpretation of intuitionistic logic, BHK for Brouwer-Heyting-Kolmogorov.\gap
We will \newlec still write $\Gamma \proves_{\text{IPC}} \varphi$ to mean that $\Gamma$ proves $\phi$ in intuitionistic propositional logic. (IQC stands for intuitionistic predicate logic, for reasons we are not sure of).
\large
\begin{center}
\textbf{Rules for IPC}
\end{center}\normalsize

\begin{enumerate}
    \item ($\wedge$-I): \begin{prooftree}
        \AxiomC{$\Gamma \proves A$}
        \AxiomC{$\Gamma \proves B$}
        \BinaryInfC{$\Gamma \proves A \wedge B$}
    \end{prooftree}
    \item ($\vee$-I) has two rules: \begin{prooftree}
        \AxiomC{$\Gamma \proves A$}
        \UnaryInfC{$\Gamma \proves A \vee B$}
    \end{prooftree}
    and \begin{prooftree}
        \AxiomC{$\Gamma \proves B$}
        \UnaryInfC{$\Gamma \proves A \vee B$}
    \end{prooftree}
    \item ($\wedge$-E) has two rules: \begin{prooftree}
        \AxiomC{$\Gamma \proves A \wedge B$}
        \UnaryInfC{$\Gamma \proves A$}
    \end{prooftree}
    and \begin{prooftree}
        \AxiomC{$\Gamma \proves A \wedge B$}
        \UnaryInfC{$\Gamma \proves B$}
    \end{prooftree}
    \item ($\vee$-E) Eliminating $\vee$ is hard, for obvious reasons, but the rule we have is:
    \begin{prooftree}
        \AxiomC{$\Gamma, A \proves C$}
        \AxiomC{$\Gamma, B \proves C$}
        \AxiomC{$\Gamma \proves A \vee B$}
        \TrinaryInfC{$\Gamma \proves C$}
    \end{prooftree}
    \item ($\perp$-E) \begin{prooftree}
        \AxiomC{$\Gamma \proves \perp$}
        \UnaryInfC{$\Gamma \proves A$}
    \end{prooftree}
    \item ($\Rightarrow$-I):
    \begin{prooftree}
        \AxiomC{$\Gamma,A \proves B$}
        \UnaryInfC{$\Gamma \proves A \Rightarrow B$}
    \end{prooftree}
    \item ($\Rightarrow$-E):
    \begin{prooftree}
        \AxiomC{$\Gamma \proves A \Rightarrow B$}
        \AxiomC{$\Gamma \proves A$}
        \BinaryInfC{$\Gamma \proves B$}
    \end{prooftree}
    \item Axiom scheme: \begin{prooftree}
        \AxiomC{}
        \UnaryInfC{$\Gamma, A \proves A$}
    \end{prooftree}
\end{enumerate}
We can obtain classical propositional logic (CPC) by adding the law of the excluded middle, or double negation exclusion:
\begin{prooftree}
    \AxiomC{}
    \UnaryInfC{$\Gamma \proves A \vee \neg A$}
\end{prooftree}
which is the law of the excluded middle, or
\begin{prooftree}
    \AxiomC{$\Gamma, \neg A \proves \perp$}
    \UnaryInfC{$\Gamma \proves A$}
\end{prooftree}
which is double negation exclusion.\gap 
We'll also use the informal notation
\begin{prooftree}
    \AxiomC{$[A]$}
    \noLine 
    \UnaryInfC{$\rvdots$}
    \noLine 
    \UnaryInfC{$X$}
    \AxiomC{$[B]$}
    \noLine
    \UnaryInfC{$\rvdots$}
    \noLine
    \UnaryInfC{$Y$}
    \BinaryInfC{$C$}
\end{prooftree}
should be read as \say{if we can prove $X$ assuming $A$ and we can prove $Y$ assuming $B$, we can infer $C$ by discharging the open assumptions $A$ and $B$.}
\begin{ex*}
    We can write an instance of ($\Rightarrow$-I) as:
    \begin{prooftree}
        \AxiomC{$\Gamma, [A]$}
        \noLine
        \UnaryInfC{$\rvdots$}
        \noLine
        \UnaryInfC{B}
        \UnaryInfC{$\Gamma \proves A \Rightarrow B$}
    \end{prooftree}
    so we removed $A$ from the set of assumptions needed for the final line. This is like the counterpart to the deduction theorem in Hilbert-style logic\footnote{\say{Part II Logic and Set Theory} style logic.}
\end{ex*}
To get intuitionistic predicate logic, we add: 
\begin{enumerate}
    \item[9.] ($\exists$-I), if $t$ is a term, then 
    \begin{prooftree}
        \AxiomC{$\Gamma \proves \varphi [ x=t ] $}
        \UnaryInfC{$\Gamma \proves \exists x. \, \varphi(x)$}
    \end{prooftree}
    \item[10.] ($\forall$-I) If $x$ is not free in $\Gamma$, then
    \begin{prooftree}
        \AxiomC{$\Gamma \proves \varphi$}
        \UnaryInfC{$\Gamma \proves \forall x. \, \varphi $}
    \end{prooftree}
    \item[11.] ($\exists$-E) When $x$ is not free in $\Gamma$ or $\psi$: \begin{prooftree}
        \AxiomC{$\Gamma \proves \exists x. \, \varphi$}
        \AxiomC{$\Gamma, \varphi \proves \psi$}
        \BinaryInfC{$\Gamma \proves \psi$}
    \end{prooftree}
    \item[12.] ($\forall$-E) If $t$ is any term:
        \begin{prooftree}
            \AxiomC{$\Gamma \proves \forall x. \, \varphi$}
            \UnaryInfC{$\Gamma \proves \varphi [x=t]$}
        \end{prooftree}
\end{enumerate}
\begin{ex}
    Let's show that $\proves_{\text{IPC}}$ $A \wedge B \Rightarrow B \wedge A$.
    \begin{prooftree}
        \AxiomC{$[A \wedge B]$}
        \UnaryInfC{$A$}
        \AxiomC{$[A \wedge B]$}
        \UnaryInfC{$B$}
        \BinaryInfC{$B \wedge A$}
        \UnaryInfC{$A \wedge B \Rightarrow B \wedge A$}
    \end{prooftree}
\end{ex}
\begin{ex} The logical axioms $\varphi \Rightarrow (\psi \Rightarrow \varphi)$ and $(\varphi \Rightarrow (\psi \Rightarrow \chi)) \Rightarrow ((\varphi \Rightarrow \psi) \Rightarrow (\varphi \Rightarrow \chi))$ are intuitionistically valid. \gap For the first, we write 
\begin{prooftree}
    \AxiomC{$[\varphi]$}
    \AxiomC{$[\psi]$}
    \BinaryInfC{$\varphi$}
    \RightLabel{\scriptsize ($\Rightarrow$-I), $\psi$}
    \UnaryInfC{$\psi \Rightarrow \varphi$}
    \RightLabel{\scriptsize ($\Rightarrow$-I), $\varphi$}
    \UnaryInfC{$\varphi \Rightarrow (\psi \Rightarrow \varphi)$}
\end{prooftree}
and this shows it. For the second axiom, we have \begin{prooftree}
    \AxiomC{$[\varphi \Rightarrow (\psi \Rightarrow \chi)]$}
    \AxiomC{$[\varphi]$}
    \RightLabel{\scriptsize($\Rightarrow$-E)}
    \BinaryInfC{$\psi \rightarrow \chi$}
    \AxiomC{$[\varphi \Rightarrow \psi]$}
    \AxiomC{$[\varphi]$}
    \RightLabel{\scriptsize($\Rightarrow$-E)}
    \BinaryInfC{$\psi$}
    \RightLabel{\scriptsize($\Rightarrow$-E)}
    \BinaryInfC{$\chi$}
    \RightLabel{\scriptsize$(\varphi)$}
    \UnaryInfC{$\varphi \Rightarrow \chi$}
    \RightLabel{\scriptsize$(\varphi \Rightarrow \psi)$}
    \UnaryInfC{$(\varphi \Rightarrow \psi) \Rightarrow (\varphi \Rightarrow \chi)$}
    \RightLabel{\scriptsize$(\varphi \Rightarrow (\psi \Rightarrow \chi))$}
    \UnaryInfC{$(\varphi \Rightarrow (\psi \Rightarrow \chi)) \Rightarrow ((\varphi \Rightarrow \psi) \Rightarrow (\varphi \Rightarrow \chi))$}
\end{prooftree}
so we can prove both of these.
\end{ex}
\begin{fancylem}{}
    If $\Gamma \proves_{\text{IPC}} \varphi$, then $\Gamma, \psi \proves_{\text{IPC}} \varphi$ for any proposition $\psi$. Moreover, if $p$ is any primitive proposition, and $\psi$ is any proposition, then $\Gamma[p:= \psi] \proves_{\text{IPC}} \varphi[p:=\psi]$.
\end{fancylem}
\proof Induction over the length of the proof. \endproof
\subsection{The simply-typed $\lambda$-calculus}
For now, we assume given a set $\Pi$ of \textbf{simple types} which is generated by a grammar: \[ \Pi = \calU \, | \, \Pi \rightarrow \Pi. \]
\begin{fancydefine}{}
    Let $V$ be an infinite set of variables. The set $\Lambda_{\Pi}$ of simply-typed lambda terms is defined by the following grammar: \[ \Lambda_{\Pi} := V \, | \, \lambda V : \Pi. \, \Lambda_{\Pi} \, | \, \Lambda_{\Pi} \, \Lambda_{\Pi} \] 
    The middle term is called $\lambda$-abstraction, and the rightmost term is called $\lambda$-application.   
\end{fancydefine}
\begin{ex*}
    If $\bbZ$ is a simple type, then we could have something like: \[ \lambda x : \bbZ.\, x^{2}\] is a simply-typed lambda term.
\end{ex*}
\begin{fancydefine*}{Context}
    A \textbf{context} is a set of pairs $\{x_{1}:\tau_{1}, \dots, x_{n} : \tau_{n}\}$ where the $x_{i}$ are (distinct) variables and $\tau_{i} \in \Pi$, i.e. the $\tau$'s are types. We write $C$ for the set of all contexts.
\end{fancydefine*}
Given a context $\Gamma \in C$, we also write $\Gamma, x.\tau$ for $\Gamma \cup \{x.\tau\}$.
\begin{fancydefine*}{Domain/range of context}
    The \textbf{domain} of $\Gamma \in C$ is the set of variables that appear in $\Gamma$ (notated $\dom(\Gamma)$), and the range of $\Gamma$ is the set of types that appear in $\Gamma$. (notated $|\Gamma|$).
\end{fancydefine*}
\begin{fancydefine}{The typability relation $\Vdash$}
    We define the typability relation $\Vdash \subseteq C \times \Delta_{\Pi} \times \Pi$ by the following: \begin{enumerate}
        \item For every $\Gamma \in C$, and every variable $x$ not occurring in $\Gamma$, and every type $\tau$, we have $\Gamma, x: \tau \Vdash x:\tau$
        \item Let $\Gamma \in C$, and $x$ be a variable not occurring in $\Gamma$, and $\sigma, \tau \in \Pi$, and let $M$ be a $\lambda$-term. If $\Gamma, x : \sigma \Vdash M : \tau$, then $\Gamma \Vdash (\lambda x : \sigma.\, M) : \sigma \to \tau$. This is saying that the lambda term which takes $x$ of type $\sigma$ goes to $M$ of type $\tau$, is of type $\sigma \to \tau$.
        \item Let $\Gamma \in C$, $\sigma,\tau \in \Pi$ be types, and $M,N$ be $\lambda$-terms. If $\Gamma \Vdash M:(\sigma \to \tau)$, and $\Gamma \Vdash N : \sigma$, then $\Gamma \Vdash (M\,N): \tau$.
    \end{enumerate}
    We will refer to the $\lambda$-calculus of $\Lambda_{\Pi}$ with this typability relation as $\lambda(\to)$.
\end{fancydefine}
A variable $x$ occurring in a $\lambda$-abstraction is \textbf{bound}, and \textbf{free} otherwise. The two functions \begin{align*}
     \lambda x: \sigma. \, x & &  \lambda y: \sigma. \, y
\end{align*}
are both the identity function of variables of type $\sigma$, so these are $\alpha$-equivalent, moreover, terms which only differ by naming of the bound variables are all $\alpha$-equivalent.\gap 
A term with no free variables is \textbf{closed}.\gap
If $M$ and $N$ are $\lambda$-terms, and $x$ is a variable, then we define the substitution of $N$ for $x$ in $M$ by: \begin{itemize}
    \item $x\, [x:=N] = N$
    \item $y\, [x:=N] = y$, if $x \neq y$.
    \item $(P\, Q)[x:=N] = P[x:= N] Q[x:=N]$, if $M=PQ$. 
    \item $\lambda y: \sigma. \, P [x:=N] = \lambda y: \sigma. \, (P[x:=N])$, if $M= \lambda y: \sigma. \, P$.
\end{itemize}
\begin{fancydefine}{$\beta$-reduction}
    The $\beta$\textbf{-reduction} relation is the smallest relation $\to_{\beta}$ on the set of $\lambda$-terms, which is closed under the following rules: \begin{itemize}
        \item $(\lambda x: \sigma. \, P)Q \to_{\beta} P[x:=Q]$
        \item $\lambda$-abstraction preserves $\beta$-reduction: \gap If $P \to_{\beta}P'$, then for all variables $x$ and types $\sigma \in \Pi$, we have that $\lambda x: \sigma. \, P \to_{\beta} \lambda x: \sigma. \, P'$.
        \item $\lambda$-application also preserves $\beta$-reduction. If $P \to_{\beta} P'$ and $Z$ is a $\lambda$-term, then $PZ \to_{\beta} P'Z$ and $ZP \to_{\beta} ZP'$.
    \end{itemize} 
    We also define $\beta$-equivalence $\cong_{\beta}$ as the smallest equivalence relation containg $\to_{\beta}$.
\end{fancydefine}
\begin{ex}
    We have $(\lambda x : \bbZ. \, (\lambda y : \tau. \, x)) 2 \to_{\beta} \lambda y: \tau.\, 2$. The function taking $x$ to a function which takes $y$ to $x$, when applied to $2$ is equivalent under $\beta$-reduction to the function which takes $y$ to $2$.
\end{ex}
When applying the first rule, we say that $(\lambda x:\, \sigma. \, P) Q$ is a $\beta$-redex and that $P[x:=Q]$ is its $\beta$-contraction.\gap 
If there is no term $N$ such that $M \to_{\beta} N$, we say that $M$ is in $\beta$-normal form. We write $M \twoheadrightarrow_{\beta} N$ if $M$ $\beta$-reduces to $N$ after (potentially multiple) applications of $\beta$-reduction.\gap We also have $\eta$-reduction: $(\lambda x: \sigma. \, (Px)) \to_{\eta} P$ when $x$ is not free in $P$.\gap 
\textbf{Shorthands:} \begin{itemize}
    \item $KLM$ for $(KL)M$
    \item $\lambda x: \sigma. \, \lambda y:\tau.\, M$ for $\lambda x: \sigma. \, (\lambda y: \tau.\, M)$
    \item $\lambda x: \sigma.\, MN$ for $\lambda x: \sigma.\, (MN)$
    \item $M \lambda x: \sigma.\, N$ for $M(\lambda x : \sigma.\, N)$.
\end{itemize}
The following can all be proved by induction over the derivation:
\begin{fancylem}{Free variables lemma}
    Suppose $\Gamma \Vdash M:\sigma.$ Then \begin{enumerate}
        \item If $\Gamma \subseteq \Gamma'$, then $\Gamma' \Vdash M: \sigma$.
        \item The free variables of $M$ occur in $\Gamma'$
        \item There is $\Gamma' \subseteq \Gamma$ with $\dom(\Gamma') = \{\,\text{free variables of $M$}\,\}$ with $\Gamma' \Vdash M: \sigma$.
    \end{enumerate}
\end{fancylem}
\begin{fancylem}{Generation Lemma}\hspace{0.1cm}
    \begin{enumerate}
        \item For every variable $x, \Gamma \in C$, and $\sigma \in \Pi$, if $\Gamma \Vdash x:\sigma$, then $x:\sigma \in \Gamma$.
        \item If $\Gamma \Vdash (MN):\sigma$, then there is a type $\tau$ such that $\Gamma \Vdash M : \tau \to \sigma$, and $\Gamma \Vdash N : \tau$.
        \item If $\Gamma \Vdash (\lambda x : \tau.\, M):\sigma$, then there are types $\tau$ and $\rho$ such that $\Gamma,x : \tau \Vdash M \, : \, \rho$. and $\sigma=(\tau \to \rho)$
    \end{enumerate}
\end{fancylem}
\begin{fancylem}{Substitution Lemma}
    The typability relation respects substitution.
\end{fancylem}
\begin{fancyprop}{Subject reduction}
    If $\Gamma \Vdash M : \sigma$ and $M \to_{\beta} N$, then $\Gamma \Vdash N : \sigma$.
\end{fancyprop}
\proof Again, by induction on the derivation of $M \to_{\beta} N$ using the generation and substitution lemmas.
\begin{fancythm}{Church-Rosser for $\lambda(\to)$}
    Suppose that we know that $M$ is typable in a context $\Gamma$, with type $\sigma$, i.e. $\Gamma \Vdash M : \sigma$, if $M \to_{\beta} N_{1}$ and $M \to_{\beta} N_{2}$, then there is a $\lambda$-term $L$ such that $N_{1} \to_{\beta} L$ and $N_{2} \to_{\beta} L$. This is sometimes called the diamond property: \begin{center}
        \begin{tikzcd}
            & M \arrow[rd, "\beta"] \arrow[ld, "\beta"] &\\
            N_{1} \arrow[rd, "\beta"] & & N_{2} \arrow[ld, "\beta"]\\
            & L &
        \end{tikzcd}
    \end{center}
\end{fancythm}
\proof ES4
\begin{corr}{}
    If a simply typed $\lambda$-term admits a $\beta$-normal form, then it is unique.
\end{corr}
\begin{fancyprop}{Uniqueness of types}
    \hspace{0.1cm}
    \begin{enumerate}
        \item If $\Gamma \Vdash M : \sigma$ and $\Gamma \Vdash M : \tau$ then $\sigma = \tau$
        \item If $\Gamma \Vdash M : \sigma$ and $\Gamma \Vdash N : \tau$, and $M \equiv_{\beta} N$, then $\sigma = \tau$.
    \end{enumerate}
\end{fancyprop}
\proof \hspace{0.1cm}
\begin{enumerate}
    \item By induction on $M$.
    \item By the hypothesis, and Church-Rosser, there is a term $L$ such that $M,N \twoheadrightarrow L$.\gap By Subject Reduction, we have $\Gamma \Vdash L : \sigma$ and $\Gamma \Vdash L : \tau$. So $\sigma = \tau$ by the previous item.
\end{enumerate}
\endproof
\begin{ex}
    There is no way to give a type to the expression $\lambda x. \, x \, x$, i.e. the lambda term applying a term to itself. If $x$ has type $\tau$, then it must have type $\tau \to \sigma$, (so that you can apply $x$ to $x$), but $\tau \neq \tau \to \sigma$.
\end{ex}
We measure the complexity of a type by looking at it as a binary tree: \[ \rho = \mu \to \left[\left( (\varphi \to \psi) \to \chi \right) \to \left( (\varphi \to \psi) \to (\varphi \to \chi) \right) \right] \] becomes \begin{center}
\end{center}
\begin{fancydefine}{height}
    The height function is the map $h: \pi \to \bbN$ that maps a type variable to $0$, and a function $\sigma \to \tau$ to $1+\max(h(\sigma),h(\tau))$. \gap We can extend $h$ to redexes: if $(\lambda x : \sigma. \, P^{\tau})^{\sigma \to \tau}. \, R^{\sigma}$ is a $\beta$-redex, we define its height to be $h(\sigma \to \tau)$.
\end{fancydefine}
\begin{fancythm}{Weak normalisation}
    If $\Gamma \Vdash M : \sigma$, then there is a finite reduction path: \[(M := )M_{0} \to_{\beta} M_{1} \to_{\beta} \dots \to_{\beta} M_{n}\] where $M_{n}$ is in $\beta$-normal form.
\end{fancythm}
\proof The idea is to induct over the complexity of $M$. \gap 
First of all, we define a function \[ m: \Lambda_{\Pi} \to \bbN \times \bbN\] by $m(M) = (0,0)$ if $M$ is in $\beta$-normal form, and $m(M) = (h(M),\redex(M))$, where $h(M)$ is the greatest height of a redex in $M$, and $\redex(M)$ is the number of redexes in $M$ that have height $h(M)$. \gap We will use lexicographic induction on the pairs $m(M)$ to show that if $M$ is typable, then $M$ can be reduced to $\beta$-normal form. \gap
\item If $\Gamma \Vdash M : \sigma$, and $M$ is in $\beta$-normal form, the claim is trivial.
\item If $M$ is not in $\beta$-normal form, let $\Delta$ be the right-most redex in $M$ of maximal height $h$. By reducing $\Delta$, we may introduce copies on existing redexes, or we can create new ones. Creation of new redexes can happen in one of the following ways: 
\begin{enumerate}
    \item If $\Delta$ is of the form $(\lambda x:\, (\rho \to \mu), \dots, x P^{\rho} \dots)(\lambda y : \rho.\, Q^{\mu})^{\rho \to \mu}$, then it reduces to $\dots(\lambda y: \rho. \, Q^{\mu})^{\rho \to \mu} P^{\rho}$, inw hich case there is a new redex of height $h(\rho \to \mu)<h$. (Because the first thing is of type $\rho \to \mu \to \tau$ for some type $\tau$)
    \item We have $\Delta = (\lambda x : \tau. \, \lambda y: \rho.\, R^{\mu})P^{\tau}$ occurring in the context $\Delta^{\rho \to \tau}Q^{\rho}$. Say $\Delta$ reduces to $\lambda y : \rho.\, R_{1}^{\mu}$. Then we create a new redex $(\lambda y: \rho.\, R_{1}^{\mu})^{\rho \to \mu}Q^{\rho}$ of height $h(\rho \to \mu) < h(\tau \to \rho \to \mu) < h$.
    \item We have $\Delta = (\lambda x: (\rho \to \mu).\, x)(\lambda y:\, \rho. \, P^{\mu})$, and that occurs in context $\Delta^{\rho \to \mu}Q^{rho}$. Reduction generates the redex $(\lambda y: \rho.\, P^{\mu})Q^{\rho}$ of height $h(\rho \to \mu)<h$.
\end{enumerate}
Now $\Delta$ itself is gone, and we just showed that any created redexes have height $<h$, so won't raise the count. There's still the possibility that reduction introduces copies of existing redexes. If $\Delta = (\lambda x: \rho \, P ^{\rho}.\, P^{\rho})Q^{\tau}$ and $P$ contais more than one free occurrence of $x$, then all the redexes in $Q$ get multiplied upon reduction. But by construction, they all have height less than $h$. \gap So if $M \to_{\beta} M'$ by reducing $\Delta$, it's always the case that $m(M')< m(M)$. By the inductive hypothesis, $M'$ can be reduced to $\beta$-NF, hence so can $M$. \endproof
\begin{fancythm}{Strong normalisation}
    Let $\Gamma \Vdash M : \sigma$. Then there is no infinite reduction $M \to_{\beta} M_{1} \to_{\beta} \cdots$.
\end{fancythm}
We will work with the fragment IPC($\to$) of IPC that only has $\to$ as the connective, and with deduction rules $(\to)$-I, $(I)$-E and (Ax), the axiom scheme.
\subsection{Propositions-as-types}
If $\calL$ is a propositional language for IPC($\to$) and $P$ is its set of primitive propositions, we can generate $\lambda(\to)$ by taking the set $U$ (of primitive types) as $P$. \gap Then both types and $\calL$ are generated by the same grammar: $U | \Pi \to \Pi$.\gap We'll think of contexts as giving a set of hypotheses, and each type $\varphi$ as the \say{type of proofs of $\varphi$.}
\begin{fancyprop}{Curry-Howard for IPC($\to$)}
    Let $\Gamma$ be a context for $\lambda(\to)$ and $\varphi$ be a proposition. Then: 
    \begin{enumerate}
        \item If $\Gamma \Vdash M : \varphi$, then \[|\Gamma| := \{ \tau \in \Pi \, : \, (x : \tau) \in \Gamma \text{ for some variable $x$ } \} \proves_{\text{IPC$(\to)$}} \varphi\]
        \item If $\Gamma \proves_{\text{IPC$(\to)$}} \varphi$, then there is a simply-typed $\lambda$-term $M$ such that \[\{(x_{\tau} : \tau) \, | \, \tau \in \Gamma \} \Vdash M : \varphi\]
    \end{enumerate}
\end{fancyprop}
\proof \hspace{0.1cm}
\begin{enumerate}
    \item We induct over the derivation of $\Gamma \Vdash M : \varphi$. 
    If $x$ is a variable not occurring in $\Gamma'$ and the derivation is of the form $\Gamma', x : \varphi \Vdash x : \varphi$, then we need to prove that $|\Gamma', x:\varphi| \proves \varphi$. 
    But this is true by the axiom scheme, since $\varphi \proves_{\text{IPC($\to$)}} \varphi$.\gap 
    If the derivation has $M$ of the form $\lambda x : \sigma.\, N$ and $\varphi = \sigma \to \tau$, then we must have that $\Gamma, x: \sigma \Vdash N: \tau$. By the induction hypothesis, we have that $|\Gamma, x: \sigma| \proves \tau$, i.e. $|\Gamma|, \sigma \proves \tau$, so $|\Gamma| \proves \sigma \to \tau = \varphi$ by arrow introduction. \gap 
    If the derivation is of the form $\Gamma \Vdash (P Q) : \varphi$, then we must have $\Gamma \Vdash P: \sigma \to \varphi$ and $\Gamma \Vdash Q : \sigma$ by the inductive hypothesis, we have that $\Gamma \proves \sigma \to \varphi$ and $|\Gamma| \proves \sigma$, so $|\Gamma| \proves \varphi$ by $(\to)$-E.
    \item Again, we induct over the derivation of $\Gamma \proves \varphi$. Write \[ \Delta := \{ x_{\tau} : \tau | \tau \in \Gamma\}.\] Say the derivation is of the form $\Gamma, \varphi \proves \varphi$ (so using the axiom scheme). If $\varphi \in \Gamma$, then clearly $\Delta \Vdash x_{\varphi} : \varphi$. If $\varphi \notin \Gamma$, then $\Delta, x_{\varphi} : \varphi \Vdash x_{\varphi} : \varphi.$\gap 
    Suppose the derivation is at a stage of the form \begin{prooftree}
        \AxiomC{$\Gamma \proves \varphi \to \psi$}
        \AxiomC{$\Gamma \proves \varphi$}
        \BinaryInfC{$\Gamma \proves \psi$}
    \end{prooftree}
    By the inductive hypothesis, there are $M,N \in \Lambda_{\Pi}$ such that $\Delta \Vdash M: (\varphi \to \psi)$ and $\Delta \Vdash N : \varphi$, so we deduce that $\Delta \Vdash (M N) : \psi$.\gap 
    Finally, if the derivation at a given stage is applying $(\to)$-I, i.e. of the form \begin{prooftree}
        \AxiomC{$\Gamma, \varphi \proves \psi$}
        \UnaryInfC{$\Gamma \proves \varphi \to \psi$}
    \end{prooftree}
    Then we have two cases. First, if $\varphi \in \Gamma$, then the inductive hypothesis gives $\Delta \Vdash M : \psi$ for some $M \in \Lambda_{\Pi}$. By weakening, we have that $\Delta, x:\varphi \Vdash M : \psi$, where $x$ is a variable not in $\Delta$. But then $\Delta \Vdash (\lambda x : \varphi.\, M) : (\varphi \to \psi)$ as needed. \gap 
    Second, if $\varphi \notin \Gamma$, then the induction hypothesis gives $\Delta, x_{\varphi} : \varphi \Vdash M: \psi$ for some $M \in \Lambda_{\Pi}$, thus $\Delta \Vdash (\lambda x_{\varphi} : \varphi.\, M): \varphi \to \psi$ as needed. 
\end{enumerate}
\endproof
\begin{ex}
    Let $\varphi$, $\psi$ be primitive propositions, and consider the term: \[ \lambda_{f} : (\varphi \to \psi) \to \varphi.\, \lambda g : \varphi \to \psi.\, g(fg).\]
    The entire term has type $((\varphi \to \psi) \to \varphi) \to ((\varphi \to \psi) \to \psi)$, so it encodes a proof of this proposition in IPC$(\to)$.\gap So I have a hypothesis $g: \varphi \to \psi$, and a hypothesis $f: (\varphi \to \psi) \to \varphi$. Then we first apply $f$ to $g$ and get a term of type $\varphi$. \begin{prooftree}
        \AxiomC{$g : [\varphi \to \psi]$}
        \AxiomC{$f: [(\varphi \to \psi) \to \varphi]$}
        \BinaryInfC{$(fg) : \varphi$}
        \AxiomC{$g : [\varphi \to \psi]$}
        \BinaryInfC{$g(fg) : \psi$}
        \UnaryInfC{$\lambda g.\, g(fg) : (\varphi \to \psi) \to \psi$} 
        \UnaryInfC{$\lambda f.\, \lambda g.\, g(fg) : ((\varphi \to \psi) \to \varphi) \to ((\varphi \to \psi) \to \varphi)$}
    \end{prooftree}
\end{ex}
\begin{fancydefine}{Full simply-typed $\lambda$-calculus}
    The types of ST$\lambda$C are generated by the following grammar: \[ \Pi = U \,|\, \Pi \to \Pi \,| \,\Pi \times \Pi\, | \,\Pi + \Pi\, |\, 0 \,|\, 1.\]
    The terms are given by the grammar: \[ V \, | \, \lambda V : \, \Pi .\, \Lambda_{\Pi} \, | \, \Lambda_{\Pi} \Lambda_{\Pi} \, | \, \langle \Lambda_{\Pi}, \Lambda_{\Pi} \rangle \, | \, \pi_{1}(\Lambda_{\Pi}) \, | \, \pi_{2}(\Lambda_{\Pi}) \, | \, \iota_{1}(\Lambda_{\Pi}) \, | \, \iota_{2}(\Lambda_{\Pi}) \, | \] \[\, \text{case}(\Lambda_{\Pi}; V. \, \Lambda_{\Pi}; V.\, \Lambda_{\Pi}) \, | \, * \, | \, !_{\Pi}\Lambda_{\Pi}\]
\end{fancydefine}
where $V$ is an infinite set of variables and $*$ is a \newlec constant.
This comes with new typing rules, first for product types: \begin{enumerate}
    \item \begin{prooftree}
    \AxiomC{$\Gamma \Vdash \psi \times \varphi$}
    \UnaryInfC{$\Gamma \Vdash \pi_{1}(M) : \psi$}
\end{prooftree}
    \item \begin{prooftree}
        \AxiomC{$\Gamma \Vdash \psi \times \varphi$}
    \UnaryInfC{$\Gamma \Vdash \pi_{2}(M) : \psi$}
    \end{prooftree}
\item \begin{prooftree}
    \AxiomC{$\Gamma \Vdash M : \psi$}
    \AxiomC{$\Gamma \Vdash M : \varphi$}
    \BinaryInfC{$\Gamma \Vdash \langle M, N \rangle : \psi \times \varphi$}
\end{prooftree}
\end{enumerate}
Then, dually, for coproduct types, we have:\begin{enumerate}
    \item \begin{prooftree}
    \AxiomC{$\Gamma \Vdash M : \psi$}
    \UnaryInfC{$\Gamma \Vdash \iota_{1}(M) : \psi + \varphi$}
\end{prooftree}
\item \begin{prooftree}
    \AxiomC{$\Gamma \Vdash N : \varphi$}
    \UnaryInfC{$\Gamma \Vdash \iota_{2}(N) : \psi + \varphi$}
\end{prooftree}
\end{enumerate}
Then for the case operator we have \begin{prooftree}
    \AxiomC{$\Gamma \Vdash L : \psi + \varphi$}
    \AxiomC{$\Gamma, x : \psi \Vdash M : \rho$}
    \AxiomC{$\Gamma, y : \varphi: \Vdash N : \rho$}
    \TrinaryInfC{$\Gamma \Vdash \textbf{case}(L;x^{\psi} .\, M;y^{\varphi} .\, N) : \rho$}
\end{prooftree}
This stands for \say{in case $L$ is something of type $\varphi$, I return $M$ with $x$ replaced by $L$, otherwise I return $N$ with $y$ replaced by $L$.}
Then also we have \begin{prooftree}
    \AxiomC{}
    \UnaryInfC{$\Gamma \Vdash * : 1$}
\end{prooftree}
\begin{prooftree}
    \AxiomC{$\Gamma \Vdash M : 0$}
    \UnaryInfC{$\Gamma \Vdash !_{\varphi} M : \varphi$}
\end{prooftree}
The first thing says that we can prove our constant is true, and that from false (i.e. $0$) we can prove anything.\gap 
This typing relation captures the BHK interpretation when paired with the new reduction rules. First, for products \begin{enumerate}
    \item $\pi_{1}(\langle M, N \rangle) \to_{\beta} M$, and 
    \item $\pi_{2}(\langle M, N \rangle) \to_{\beta} N$.
    \item $\langle \pi_{1}, \pi_{2} \rangle \to_{\eta} M$.
\end{enumerate}
Then for cases we have: \begin{enumerate}
    \item case$(\iota_{1};x.\, K, y.\, L) \to_{\beta} K[x:=M]$, and 
    \item case$(\iota_{2};x.\, K, y.\, L) \to_{\beta} K[y:=M]$
\end{enumerate}
And then for our constant $*$, we have \begin{itemize}
    \item If $\Gamma \Vdash M : 1$, then $M \to_{\eta} *$.
\end{itemize}
We let $0$ correspond to $\perp$, product types to conjunctions, coproduct types to disjunctions, to once again have propositions as types.\gap 
Redexes are now the expressions consisting of a constructor (pair formation, $\lambda$-abstraction, injections) followed by the corresponding destructor (projections, applications, or case-expressions).
\begin{ex}
    Consider the following proof of $\varphi \wedge \chi \to (\psi \to \varphi)$. \begin{prooftree}
        \AxiomC{$[\varphi \to \chi]$}
        \RightLabel{\scriptsize($\wedge$-E)}
        \UnaryInfC{$\chi$}
        \AxiomC{$[\psi]$}
        \RightLabel{\scriptsize($\Rightarrow$-I)}
        \BinaryInfC{$\psi \to \varphi$}
        \RightLabel{\scriptsize($\Rightarrow$-I)}
        \UnaryInfC{$\varphi \wedge \chi \to (\psi \to \varphi)$}
    \end{prooftree}
This corresponds to the $\lambda$-term which is given by $\lambda p :\, \varphi \times \chi.\, \lambda b: \, \psi.\, \pi_{1}(p)$ (can work through it by using projections for elimination, $\lambda$-abstraction for ($\Rightarrow$-I), and considering the fact that $\varphi \wedge \chi$ is of type $\varphi \times \chi$).\end{ex}

With these rules, we can extend Curry-Howard to the whole of IPC, and ST$\lambda$C. Where we have \begin{align*}
    \text{ST$\lambda$C} & &\text{IPC}\\
    \text{(prim) types} & &\text{(prim) propositions}\\
    \text{variables} & &\text{hypotheses}\\
    \text{$\lambda$-term} & &\text{proof}\\
    \text{term inhabitation} & &\text{provability}\\
    \text{term reduction} & &\text{proof normalisation}\\
\end{align*}

\subsection{Semantics of IPC}
\begin{fancydefine}{Heyting Algebra}
    A \textbf{Heyting algebra} is a bounded lattice eequipped with a binary operation $\Rightarrow : H \times H \to H$ such that \[ A \wedge B \leq C \,\,\, \text{iff} \,\,\, A \leq (B \Rightarrow C) \,\,\,\,\, \forall A,B,C \in H.\]

    A morphism of Heyting algebras is a function that preserves all finite meets (including $1$, the $0$-ary meet), finite joins (including $0$), and $\Rightarrow$. 
\end{fancydefine}

\begin{exs}\hspace{0.1cm}
    \begin{enumerate}
        \item Every Boolean algebra is a Heyting algebra, where $a \Rightarrow b := \neg a \vee b$. Note that $\neg a := a \to \perp$.
        \item Every topology on a set $X$ is a Heyting algebra $U \Rightarrow V := \Int\left( (X \setminus U) \cup V)\right)$
        \item It turns out that every finite distributive lattice is a Heyting algebra (see ES4).
        \item Your mum is a Heyting algebra.
        \item The Lindenbaum-Tarski algebra of a propositional theory $\calT$ with respect to IPC is a Heyting algebra (exercise).
    \end{enumerate}
\end{exs}
\begin{fancydefine}{$H$-valuation}
    Let $H$ be a Heyting algebra and $\calL$ be a propositional language with a set $P$ of primitive propositions.
    An $H$-\textbf{valuation} is a function \[ v : P \to H \] extended to the whole of $\calL$ by setting: \begin{itemize}
        \item $v(\perp) := \perp$
        \item $v(A \wedge B) := v(A) \wedge v(B)$
        \item $v(A \vee B) := v(A) \vee v(B)$
        \item $v(A \to B) := v(A) \Rightarrow v(B)$
    \end{itemize}
    A proposition $A$ is $H$-valid if $v(A) = \top$ for all $H$-valuations $v$, and is an $H$-consequence of a finite set of propositions if $v(\bigwedge \Gamma) \leq v(A)$ (written $\Gamma \vDash_{H} A$).
\end{fancydefine}
\begin{fancylem}{Soundness of the Heyting semantics}
    Let $H$ be a Heyting algebra and $v : \calL \to H$ be a valuation, then \[ \Gamma \proves_{\text{IPC}} A \Rightarrow \Gamma \models_{H} A. \]
\end{fancylem}
\proof By induction over the derivations of $\Gamma \proves A$.\gap There are only a few ways you can prove things: \begin{enumerate}
    \item[(Ax)] The proposition holds, as \[v( (\bigwedge \Gamma) \wedge A) = v\left(\bigwedge \Gamma\right) \wedge v(A) \leq v(A).\]
    \item[($\wedge$-I)] Then $A = B \wedge C$ and we have derivations $\Gamma_{1} \proves B$ and $\Gamma_{2} \proves C$ with $\Gamma_{1},\Gamma_{2} \subseteq \Gamma$. Then by the induction hypothesis, we know that $v\left(\bigwedge \Gamma_{1}\right) \leq v(B)$ and $v\left( \bigwedge \Gamma_{2} \right) \leq v(C)$. Thus \begin{align*} 
        v(\Gamma) &\leq v(\Gamma_{1}) \wedge v(\Gamma_{2})\\
                  &\leq v(B) \wedge v(C)\\ 
                  &\leq v(B \wedge C) = v(A)
    \end{align*}
    \item[($\to$)-I] Then $A = B \to C$, and therefore we must have $\Gamma \cup \{B\} \proves C$. By the inductive hypothesis, we have that $v\left(\bigwedge \Gamma \right) \wedge v(B) = v(\bigwedge \Gamma \wedge B) \leq v(C)$. Thus \begin{align*} 
        v\left(\bigwedge \Gamma\right) &\leq (v(B) \Rightarrow v(C))\\
                            = v(B \to C)
    \end{align*}
    by definition of \say{$\Rightarrow$}
    \item[($\vee$)-I] Then $A = B \vee C$, and w.l.o.g. we have $\Gamma \proves B$. By the inductive hypothesis, we have $v(\bigwedge \Gamma) \leq v(B)$. But $v(B \vee C) = v(B) \vee v(C)$ and $v(B) \leq v(B) \vee v(C)$. Therefore $v(\Gamma) \leq v(A)$.
    \item[($\wedge$)-E] By the inductive hypothesis, we have \begin{align*}
        v\left( \bigwedge \Gamma\ \right) &\leq v(A \wedge B)\\
        &= v(A) \wedge v(B)\\
        &\leq v(A),v(B).
    \end{align*}
    \item[($\to$)-E] We know that $v(A \to B) = v(A) \Rightarrow v(B)$. From $v(A \to B) \leq (v(A) \Rightarrow v(B))$, we deduce that $v(A) \wedge v(A \to B) \leq v(B)$.\gap So if $v\left( \bigwedge \Gamma \right) \leq v(A \to B)$ and $v\left( \bigwedge \Gamma \right) \leq v(A)$ then $V\left( \bigwedge \Gamma \right) \leq v(B)$ as needed.
    \item[($\vee$-E)] By the inductive hypothesis $v\left(A \wedge \bigwedge \Gamma \right) \leq v(C)$ and $v\left( B \wedge \bigwedge \Gamma \right) \leq v(C)$, and also $v(\bigwedge \Gamma) \leq v(A \vee B) = v(A) \vee v(B)$. This last fact means that \[ v\left( \bigwedge \Gamma \right) = v\left(\bigwedge \Gamma \right) \wedge \left( v(A) \vee v(B) \right),\]
    which is the same as \[\left( v\left( \bigwedge \Gamma \right) \vee v(A) \right) \vee \left( v\left( \bigwedge \Gamma \right) \vee v(B) \right)\] since Heyting algebras are distributive lattices (c.f. ES4). This last term is bounded by $v(C)$, by the initial inequalities.
    \item[($\perp$-E)] If $v\left( \bigwedge \Gamma \right) \leq v(\perp) = \perp$, then $v\left( \bigwedge \Gamma \right) = \perp$, in which case $v\left( \bigwedge \Gamma \right) \leq v(A)$ for any $A$ by minimality.
\end{enumerate}
\endproof 
\begin{ex}
    The LEM is \textit{not} provable in IPC: let $p$ be a primitive proposition, and consider the Heyting algebra given by the topology $\{ \emptyset, \{1\}, \{1,2\}\}$ on $\{1,2\}$. We can define a valuation $v$ with $v(p) = \{1\}$, in which case $v(\neg p) = \neg \{ p\} = \Int(\{1,2\}\setminus \{1\}) = \emptyset$, and so $v(p \vee \neg p) = \{ 1\} \neq \top$. \gap So soundness iplies that $p \vee \neg p$ is not provable in IPC.
\end{ex}
\begin{ex}
    Pierce's law, $\left(\left( p \to q \right) \to p \right) \to p$, is not intuitionistically valid. Take $H$ to be the usual topology on $\bbR^{2}$ and our valuation $v$ takes \begin{align*} p &\mapsto \bbR^{2} \setminus \{\mathbf{0}\} \\ q &\mapsto \emptyset \end{align*}
\end{ex}
Recall that classical completeness says $\Gamma \proves_{\text{CPC}} A$ iff $\Gamma \models_{2} A$ where $2$ is the Boolean algebra $2$. For IPC we can't replace $2$ with any finite Heyting algebra (see ES4), but we can still prove the following: \begin{fancythm}{Completeness of Heyting semantics}
    A proposition is provable in IPC if and only if it is $H$-valid for every Heyting algebra $H$.
\end{fancythm}
\proof \hspace{0.1cm} \begin{enumerate}
    \item[$(\Rightarrow)$] If $\proves_{\text{IPC}}A$ then $\top = v(\wedge \emptyset) \leq v(A)$ for any Heyting algebra $H$ and valuation $v$, by soundness. Then $v(A) = \top$, and thus $A$ is $H$-valid.
    \item[$(\Leftarrow)$] Recall that the Lindebaum-Tarski algebra: $\calL / \sim$ for the empty theory with respect to IPC is a Heyting algebra. Define a valuation $v$ by extending the map from the set of propositins $\bbP$ \begin{align*}
        \bbP &\to \calL/ \sim \\
        p \to [p]\\
    \end{align*}
    to all propositions in IPC. By induction, we can show that $v(\varphi) = [\varphi]$ for every proposition $\varphi \in \calL$. Now, $A$ is valid in every Heyting algebra, and w.r.t. any valuation, so in particular $v(A) = \top$ in $\calL / \sim$. But that means that $\proves_{\text{IPC}} A \leftrightarrow \top$, i.e. IPC proves $A$.
\end{enumerate}
\endproof
Given a poset $S$, we can construct sets $a\uparrow := \{ s \in S : a \leq s\}$, called principal up-sets (everything greater than some element). Recall that $\calU \subseteq S$ is a terminal set iff $a \uparrow \subseteq \calU$ for all $a \in \calU$.
\begin{fancyprop}{}
    If $S$ is a poset, then the set $T(S)$ of terminal segments of $S$ is a Heyting algebra.
\end{fancyprop}
\proof Meets and joins are intersections and unions, and the order is by inclusion. The implication is given by \[ U \Rightarrow V = \{ s \in S: \, (s \uparrow) \cap U \subseteq V\}. \]
Then $W \subseteq U \Rightarrow V$ iff $U \cap W \subseteq V$.
\begin{fancydefine}{Kripke Model}
    Let $P$ be a set of primitive propositions. A Kripke model is a triple $(S, \leq, \Vdash)$, where $S$ is a poset and $\Vdash \subseteq S \times P$ (which we call forcing) is a relation satisfying the persistence property (i.e. persistence between worlds). If $p \in P$ is such that $s \Vdash p$ and $s \leq s'$ then $s' \Vdash p$.
\end{fancydefine} 
Every valuation $v$ on $T(S)$\footnote{The terminal segments of $S$, a Heyting algebra} induces a Kripke model by setting that $s \Vdash p$ iff $s \in v(p)$.
\begin{fancydefine}{}
    Let $(S, \leq, \Vdash)$ be a Kripke model. We can extend the forcing relation to a relation $\Vdash \subseteq S \times \calL$ recursively as follows: \begin{itemize}
        \item There's no $s \in S$ with $s \Vdash \perp$,
        \item $s \Vdash \varphi \wedge \psi$ iff $s \Vdash \varphi$ and $s \Vdash \psi$,
        \item $s \Vdash \varphi \vee \psi$ iff $s \Vdash \varphi$ or $s \Vdash \psi$,
        \item $s \Vdash \varphi \Rightarrow \psi$ iff $s' \Vdash \varphi \Rightarrow s' \Vdash \psi$ for every $s' \geq s$.
    \end{itemize}
\end{fancydefine}
We can check that persistence holds for arbitrary propositions. Moreover, \begin{itemize}
    \item $s \Vdash \neg \varphi$ iff $s' \nVdash \varphi$ for all $s' \geq s$,
    \item $s \Vdash \neg \neg \varphi$ iff for each $s' \geq s$, there is $s'' \geq s'$ with $s'' \Vdash \varphi$.
\end{itemize}
We say that $S \Vdash \varphi$ if all the worlds force it (if $S$ has a least element, this is equivalent to saying that the least element forces $\varphi$).
\begin{exs}
    Consider the following Kripke models \begin{enumerate}
        \item \hspace{0.1cm}
        \begin{center}
            \begin{tikzcd}
                s' \Vdash p \arrow[d,dash]\\
                s
            \end{tikzcd}
        \end{center}
        \item \hspace{0.1cm}
        \begin{center}
            \begin{tikzcd}
                s' \arrow[dr, dash] & & s'' \Vdash p \arrow[dl, dash]\\
                & s &
            \end{tikzcd}
        \end{center}
        \item \hspace{0.1cm}
        \begin{center}
            \begin{tikzcd}
                s' \Vdash p, \Vdash q \arrow[dash,d]\\
                s
            \end{tikzcd}
        \end{center}
    \end{enumerate}
    Note that: In $1$, we have $s \nVdash \neg p$ since $s' \geq s$ and $s' \Vdash p$. We also know that $s \nVdash p$. Thus $s \nVdash p \vee \neg p$. It is also the case that $s \Vdash \neg \neg p$, yet $s \nVdash p$, so $s \nVdash \neg \neg p \rightarrow p$ either. \gap 
    In $2$, we have $s \nVdash \neg \neg p$, since there is a world $s' \geq s$ which can't access a world that forces $p$. But $s \nVdash \neg p$, since $s$ can see a world that does believe in $p$. So $s \nVdash \neg \neg p \vee \neg p$. \gap 
    In $3$, $s \nVdash (p \rightarrow q) \rightarrow (\neg p \vee q)$. All worlds force $p \rightarrow q$, and $s \nVdash q$. So to check the claim, it is enough to show that $s \nVdash \neg p$. But this is clearly the case, since $s' \geq s$ and $s' \Vdash p$.
\end{exs}
\begin{fancylem}{}
    Let $H$ be a Heyting algebra and $v$ be an $H$-valuation. Then we can build a Kripke model $(S, \leq, \Vdash)$ such that \[ 
        v \vDash_{H} \varphi \text{ iff } S \Vdash \varphi,\] for each proposition $\varphi$.
\end{fancylem}
\proof Let $S$ be the set of prime filters of $H$, ordered by $\subseteq$. We write $F \Vdash p$ iff $v(p) \in F$ for primitive propositions $p$, and we prove by induction that this extends to arbitrary propositions. \gap 
For implications, say that $F \Vdash (\psi \rightarrow \psi')$ and suppose that $v(\psi \rightarrow \psi') = (v(\psi) \Rightarrow v(\psi')) \notin F$. Let $G'$ be the least filter containing $F$ and $v(\psi)$. Then \[ G' = \{ b  \, : \, \exists f \in F.\,\, f \wedge v(\psi) \leq b \} \] and $v(\psi') \notin G'$, since otherwise $f \wedge v(\psi) \leq v(\psi')$ for some $f \in F$, so that $f \leq v(\psi \rightarrow \psi') \in F$. In particular, $G'$ is proper, so we can find a prime filter extending $G'$\footnote{this is dual to the Boolean prime ideal theorem.}. So let $G$ be such a prime filter extending $G'$ which does not contain $v(\psi')$, which exists by Zorn's lemma. \gap 
By the induction hypothesis, $G \Vdash \psi$, and since $F \Vdash \psi \rightarrow \psi'$ and $F \subseteq G' \subseteq G$, we must have $G \Vdash \psi'$. But then $v(\psi') \in G$, which is a contradiction. This shows that $F \Vdash \psi \rightarrow \psi'$ implies $v(\psi \rightarrow \psi') \in F$. \gap 
Conversely, say that $v(\psi \rightarrow \psi') \in F \subseteq G \Vdash \psi$. By the inductive hypothesis, $v(\psi) \in G$, and so $v(\psi) \Rightarrow v(\psi') \in G$ (as $F \subseteq G$). But then $v(\psi') \geq v(\psi) \wedge (v(\psi) \Rightarrow v(\psi')) \in G$, so from the inductive hypothesis we get $G \Vdash \psi'$ as needed.\gap The other connectives are easy, but primeness is needed for disjunctions. All that's left is to show that $v \Vdash_{H}$ iff $S \Vdash \varphi$. \gap 
If $v \Vdash_{H} \varphi$, then $v(\varphi)= \top$, so $\varphi$ is in every filter of $H$, so $F \Vdash \varphi$ for every prime filter. \gap 
Conversely, if $S \Vdash \varphi$ but $v \nvDash_{H} \varphi$. Then since $v(\varphi) \neq \top$, there must be a proper filter which does not contain it. Take this filter, and extend it by Zorn's lemma to a prime filter $G$ which does not contain $v(\varphi)$, i.e. $G \nVdash \varphi$. This gives a contradiction. \endproof
\begin{fancythm}{Completeness of the Kripke Semantics}
    For every proposition $\varphi$, we have that $\Gamma \proves_{\text{IPC}} \varphi$ iff for all Kripke models $(S, \leq, \Vdash)$, the condition $S \Vdash \Gamma$ implies $S \Vdash \varphi$.
\end{fancythm}
\proof
Soundness is by induction as usual. \gap 
Say that $\Gamma \nvdash_{\text{IPC}} \varphi$. Then $v \vDash_{H} \Gamma$ but $v \nvDash_{H}\varphi$ for some Heyting algebra $H$ and valuation $v$ (by theorem (2.4.7)). Then we apply lemma 2.4.12 to $H$ and $v$, and obtain a Kripke model $(S, \leq, \Vdash)$ such that $S \Vdash \Gamma$, but $S \nVdash \varphi$, contradicting the hypothesis we have on every Kripke model.
\endproof















\end{document}